\IfFileExists{../diagrams/huang_templates.tex}{\definecolor{huangPurple}{RGB}{113,49,169}
\definecolor{huangBlue}{RGB}{65,109,190}
\definecolor{huangRed}{RGB}{196,0,0}
\definecolor{huangOrange}{RGB}{239,126,45}
\definecolor{huangGold}{RGB}{196,145,15}
\definecolor{huangGreen}{RGB}{91,164,61}
\definecolor{huangGray}{RGB}{120,120,120}
\definecolor{huangPaleBlue}{RGB}{224,234,238}
\definecolor{huangPaleGreen}{RGB}{225,240,216}
\definecolor{huangPeach}{RGB}{246,185,139}
\definecolor{huangLavender}{RGB}{216,205,236}

\tikzset{
  huang arrow/.style={-Stealth, thick, draw=huangGray},
  huang dashed box/.style={draw, dashed, rounded corners=5pt, fill=huangPaleBlue, inner sep=5pt},
  huang bit/.style={draw, dashed, rounded corners=4pt, minimum width=0.42cm, minimum height=0.36cm, inner sep=1pt, font=\scriptsize\bfseries},
  huang gate/.style={draw, rounded corners=2pt, fill=huangPeach, minimum width=0.35cm, minimum height=0.54cm, inner sep=1pt},
  huang meas/.style={draw, rounded corners=2pt, fill=huangLavender, minimum width=0.24cm, minimum height=0.24cm, inner sep=0pt},
  huang label/.style={font=\scriptsize},
  huang title/.style={font=\bfseries}
}

\providecommand{\HuangAtom}[2]{%
  \begin{scope}[shift={(#1,#2)}, scale=0.18]
    \foreach \a in {0,60,...,300} {
      \draw[huangPurple, thick, rotate=\a] (0,0) ellipse (0.36 and 0.9);
    }
    \fill[huangGreen!85!black] (0,0) circle (0.11);
  \end{scope}%
}

\providecommand{\HuangLabPatch}[3][1]{%
  \begin{scope}[shift={(#2,#3)}, scale=#1]
    \node[huang dashed box, minimum width=2.25cm, minimum height=1.55cm] (labpatch) at (0,0) {};
    \foreach \x/\y in {-0.7/0.43,0/0.43,0.7/0.43,-0.7/-0.05,0/-0.05,0.7/-0.05,-0.35/-0.53,0.35/-0.53}
      {\HuangAtom{\x}{\y}}
  \end{scope}%
}

\providecommand{\HuangLabStack}[3][1]{%
  \begin{scope}[shift={(#2,#3)}, scale=#1]
    \foreach \dx in {-0.42,-0.18,0.06} {
      \node[huang dashed box, minimum width=1.75cm, minimum height=1.25cm, fill=huangPaleBlue!85] at (\dx,0) {};
    }
    \foreach \x/\y in {-0.58/0.36,0/0.36,0.58/0.36,-0.58/-0.03,0/-0.03,0.58/-0.03,-0.3/-0.43,0.3/-0.43}
      {\HuangAtom{\x}{\y}}
  \end{scope}%
}

\providecommand{\HuangMonitor}[3][1]{%
  \begin{scope}[shift={(#2,#3)}, scale=#1]
    \fill[black] (-0.72,-0.54) rectangle (0.72,0.55);
    \fill[huangLavender!60] (-0.6,-0.4) rectangle (0.6,0.4);
    \fill[white, opacity=0.72] (-0.55,-0.35) -- (-0.12,0.35) -- (0.6,0.35) -- (0.6,-0.35) -- cycle;
    \node[font=\scriptsize\bfseries, text=huangGray!80] at (0.12,0.02) {Model $\Psi$};
    \fill[black] (0,-0.47) circle (0.06);
    \fill[black] (-0.13,-0.8) rectangle (0.13,-0.54);
    \draw[black, line width=1.5pt] (-0.38,-0.84) -- (0.38,-0.84);
  \end{scope}%
}

\providecommand{\HuangMeasurementStack}[3][1]{%
  \begin{scope}[shift={(#2,#3)}, scale=#1]
    \foreach \i in {0,...,6} {
      \node[huang meas] at (0,0.18*\i) {};
      \draw[huangGold!80!black] (-0.08,0.18*\i-0.02) -- (0.08,0.18*\i+0.07);
    }
  \end{scope}%
}

\providecommand{\HuangCircuitGrid}[4][1]{%
  \begin{scope}[shift={(#2,#3)}, scale=#1]
    \foreach \x in {-0.75,-0.45,-0.15,0.15,0.45,0.75}
      \draw[huangGray] (\x,0.9) -- (\x,-1.05);
    \foreach \x/\y/\w/\h in {
      -0.75/0.55/.34/.45,-0.45/0.18/.24/.42,-0.15/0.55/.24/.52,0.15/0.18/.50/.42,0.45/0.55/.34/.45,0.75/0.18/.24/.42,
      -0.75/-0.25/.24/.46,-0.45/-0.63/.52/.42,-0.15/-0.25/.24/.52,0.15/-0.63/.52/.42,0.45/-0.25/.24/.46,0.75/-0.63/.34/.42}
      \node[huang gate, minimum width=\w cm, minimum height=\h cm] at (\x,\y) {};
    \node[font=\scriptsize] at (0,1.12) {$#4$};
  \end{scope}%
}

\providecommand{\HuangNeuralNet}[3][1]{%
  \begin{scope}[shift={(#2,#3)}, scale=#1]
    \node[draw, dashed, rounded corners=8pt, fill=huangPaleGreen, minimum width=1.65cm, minimum height=2.45cm] at (0,0) {};
    \foreach \x in {-0.45,0,0.45} \fill[huangBlue!38] (\x,0.75) circle (0.08);
    \foreach \x in {-0.58,-0.18,0.22,0.58} \fill[huangBlue!38] (\x,0.15) circle (0.08);
    \foreach \x in {-0.58,-0.18,0.22,0.58} \fill[huangBlue!38] (\x,-0.45) circle (0.08);
    \fill[huangBlue!38] (0,-1.02) circle (0.08);
    \foreach \a in {-0.45,0,0.45} \foreach \b in {-0.58,-0.18,0.22,0.58}
      \draw[huangGray!65, thin] (\a,0.75) -- (\b,0.15);
    \foreach \a in {-0.58,-0.18,0.22,0.58} \foreach \b in {-0.58,-0.18,0.22,0.58}
      \draw[huangGray!65, thin] (\a,0.15) -- (\b,-0.45);
    \foreach \a in {-0.58,-0.18,0.22,0.58}
      \draw[huangGray!65, thin] (\a,-0.45) -- (0,-1.02);
  \end{scope}%
}

\providecommand{\HuangDeviceIcon}[3][1]{%
  \begin{scope}[shift={(#2,#3)}, scale=#1]
    \fill[brown!80!black] (-0.72,0.82) rectangle (0.72,1.12);
    \fill[brown!70!black] (-0.82,0.72) rectangle (0.82,0.78);
    \foreach \x in {-0.48,-0.28,0.28,0.48}
      \draw[line width=1.5pt, huangGold] (\x,0.7) -- (\x,-0.55);
    \foreach \y in {0.42,0.12,-0.18}
      \draw[line width=1.2pt, brown!70!black] (-0.65,\y) -- (0.65,\y);
    \foreach \x in {-0.48,-0.38,0.38,0.48}
      \foreach \y in {0.28,-0.08}
        \draw[line width=1pt, huangGold] (\x,\y) circle (0.055);
    \foreach \x in {-0.52,-0.42,0.42,0.52}
      \foreach \y in {-0.38,-0.48,-0.58}
        \draw[line width=1.5pt, huangGold] (\x-0.08,\y) -- (\x+0.08,\y);
    \fill[huangGold!70] (-0.28,-0.6) rectangle (0.28,-0.85);
    \fill[brown!85!black] (-0.15,-0.85) rectangle (0.15,-1.18);
  \end{scope}%
}

\providecommand{\HuangHypercube}[3][1]{%
  \begin{scope}[shift={(#2,#3)}, scale=#1]
    \coordinate (z000) at (-1.2,-0.25);
    \coordinate (z001) at (-0.25,0.85);
    \coordinate (z010) at (-0.25,-0.25);
    \coordinate (z011) at (0.65,0.85);
    \coordinate (z100) at (-0.25,-1.2);
    \coordinate (z101) at (0.65,-0.25);
    \coordinate (z110) at (0.65,-1.2);
    \coordinate (z111) at (1.55,-0.25);
    \foreach \a/\b in {z000/z001,z000/z010,z000/z100,z001/z011,z001/z101,z010/z011,z010/z110,z100/z101,z100/z110,z011/z111,z101/z111,z110/z111}
      \draw[line width=1.5pt, huangGray] (\a) -- (\b);
    \foreach \p/\lab/\anch in {z000/000/below left,z001/001/above,z010/010/above left,z011/011/above,z100/100/below,z101/101/above,z110/110/below,z111/111/right}
      \fill[huangRed] (\p) circle (0.08) node[\anch, font=\scriptsize, text=black] {\lab};
  \end{scope}%
}
}{\definecolor{huangPurple}{RGB}{113,49,169}
\definecolor{huangBlue}{RGB}{65,109,190}
\definecolor{huangRed}{RGB}{196,0,0}
\definecolor{huangOrange}{RGB}{239,126,45}
\definecolor{huangGold}{RGB}{196,145,15}
\definecolor{huangGreen}{RGB}{91,164,61}
\definecolor{huangGray}{RGB}{120,120,120}
\definecolor{huangPaleBlue}{RGB}{224,234,238}
\definecolor{huangPaleGreen}{RGB}{225,240,216}
\definecolor{huangPeach}{RGB}{246,185,139}
\definecolor{huangLavender}{RGB}{216,205,236}

\tikzset{
  huang arrow/.style={-Stealth, thick, draw=huangGray},
  huang dashed box/.style={draw, dashed, rounded corners=5pt, fill=huangPaleBlue, inner sep=5pt},
  huang bit/.style={draw, dashed, rounded corners=4pt, minimum width=0.42cm, minimum height=0.36cm, inner sep=1pt, font=\scriptsize\bfseries},
  huang gate/.style={draw, rounded corners=2pt, fill=huangPeach, minimum width=0.35cm, minimum height=0.54cm, inner sep=1pt},
  huang meas/.style={draw, rounded corners=2pt, fill=huangLavender, minimum width=0.24cm, minimum height=0.24cm, inner sep=0pt},
  huang label/.style={font=\scriptsize},
  huang title/.style={font=\bfseries}
}

\providecommand{\HuangAtom}[2]{%
  \begin{scope}[shift={(#1,#2)}, scale=0.18]
    \foreach \a in {0,60,...,300} {
      \draw[huangPurple, thick, rotate=\a] (0,0) ellipse (0.36 and 0.9);
    }
    \fill[huangGreen!85!black] (0,0) circle (0.11);
  \end{scope}%
}

\providecommand{\HuangLabPatch}[3][1]{%
  \begin{scope}[shift={(#2,#3)}, scale=#1]
    \node[huang dashed box, minimum width=2.25cm, minimum height=1.55cm] (labpatch) at (0,0) {};
    \foreach \x/\y in {-0.7/0.43,0/0.43,0.7/0.43,-0.7/-0.05,0/-0.05,0.7/-0.05,-0.35/-0.53,0.35/-0.53}
      {\HuangAtom{\x}{\y}}
  \end{scope}%
}

\providecommand{\HuangLabStack}[3][1]{%
  \begin{scope}[shift={(#2,#3)}, scale=#1]
    \foreach \dx in {-0.42,-0.18,0.06} {
      \node[huang dashed box, minimum width=1.75cm, minimum height=1.25cm, fill=huangPaleBlue!85] at (\dx,0) {};
    }
    \foreach \x/\y in {-0.58/0.36,0/0.36,0.58/0.36,-0.58/-0.03,0/-0.03,0.58/-0.03,-0.3/-0.43,0.3/-0.43}
      {\HuangAtom{\x}{\y}}
  \end{scope}%
}

\providecommand{\HuangMonitor}[3][1]{%
  \begin{scope}[shift={(#2,#3)}, scale=#1]
    \fill[black] (-0.72,-0.54) rectangle (0.72,0.55);
    \fill[huangLavender!60] (-0.6,-0.4) rectangle (0.6,0.4);
    \fill[white, opacity=0.72] (-0.55,-0.35) -- (-0.12,0.35) -- (0.6,0.35) -- (0.6,-0.35) -- cycle;
    \node[font=\scriptsize\bfseries, text=huangGray!80] at (0.12,0.02) {Model $\Psi$};
    \fill[black] (0,-0.47) circle (0.06);
    \fill[black] (-0.13,-0.8) rectangle (0.13,-0.54);
    \draw[black, line width=1.5pt] (-0.38,-0.84) -- (0.38,-0.84);
  \end{scope}%
}

\providecommand{\HuangMeasurementStack}[3][1]{%
  \begin{scope}[shift={(#2,#3)}, scale=#1]
    \foreach \i in {0,...,6} {
      \node[huang meas] at (0,0.18*\i) {};
      \draw[huangGold!80!black] (-0.08,0.18*\i-0.02) -- (0.08,0.18*\i+0.07);
    }
  \end{scope}%
}

\providecommand{\HuangCircuitGrid}[4][1]{%
  \begin{scope}[shift={(#2,#3)}, scale=#1]
    \foreach \x in {-0.75,-0.45,-0.15,0.15,0.45,0.75}
      \draw[huangGray] (\x,0.9) -- (\x,-1.05);
    \foreach \x/\y/\w/\h in {
      -0.75/0.55/.34/.45,-0.45/0.18/.24/.42,-0.15/0.55/.24/.52,0.15/0.18/.50/.42,0.45/0.55/.34/.45,0.75/0.18/.24/.42,
      -0.75/-0.25/.24/.46,-0.45/-0.63/.52/.42,-0.15/-0.25/.24/.52,0.15/-0.63/.52/.42,0.45/-0.25/.24/.46,0.75/-0.63/.34/.42}
      \node[huang gate, minimum width=\w cm, minimum height=\h cm] at (\x,\y) {};
    \node[font=\scriptsize] at (0,1.12) {$#4$};
  \end{scope}%
}

\providecommand{\HuangNeuralNet}[3][1]{%
  \begin{scope}[shift={(#2,#3)}, scale=#1]
    \node[draw, dashed, rounded corners=8pt, fill=huangPaleGreen, minimum width=1.65cm, minimum height=2.45cm] at (0,0) {};
    \foreach \x in {-0.45,0,0.45} \fill[huangBlue!38] (\x,0.75) circle (0.08);
    \foreach \x in {-0.58,-0.18,0.22,0.58} \fill[huangBlue!38] (\x,0.15) circle (0.08);
    \foreach \x in {-0.58,-0.18,0.22,0.58} \fill[huangBlue!38] (\x,-0.45) circle (0.08);
    \fill[huangBlue!38] (0,-1.02) circle (0.08);
    \foreach \a in {-0.45,0,0.45} \foreach \b in {-0.58,-0.18,0.22,0.58}
      \draw[huangGray!65, thin] (\a,0.75) -- (\b,0.15);
    \foreach \a in {-0.58,-0.18,0.22,0.58} \foreach \b in {-0.58,-0.18,0.22,0.58}
      \draw[huangGray!65, thin] (\a,0.15) -- (\b,-0.45);
    \foreach \a in {-0.58,-0.18,0.22,0.58}
      \draw[huangGray!65, thin] (\a,-0.45) -- (0,-1.02);
  \end{scope}%
}

\providecommand{\HuangDeviceIcon}[3][1]{%
  \begin{scope}[shift={(#2,#3)}, scale=#1]
    \fill[brown!80!black] (-0.72,0.82) rectangle (0.72,1.12);
    \fill[brown!70!black] (-0.82,0.72) rectangle (0.82,0.78);
    \foreach \x in {-0.48,-0.28,0.28,0.48}
      \draw[line width=1.5pt, huangGold] (\x,0.7) -- (\x,-0.55);
    \foreach \y in {0.42,0.12,-0.18}
      \draw[line width=1.2pt, brown!70!black] (-0.65,\y) -- (0.65,\y);
    \foreach \x in {-0.48,-0.38,0.38,0.48}
      \foreach \y in {0.28,-0.08}
        \draw[line width=1pt, huangGold] (\x,\y) circle (0.055);
    \foreach \x in {-0.52,-0.42,0.42,0.52}
      \foreach \y in {-0.38,-0.48,-0.58}
        \draw[line width=1.5pt, huangGold] (\x-0.08,\y) -- (\x+0.08,\y);
    \fill[huangGold!70] (-0.28,-0.6) rectangle (0.28,-0.85);
    \fill[brown!85!black] (-0.15,-0.85) rectangle (0.15,-1.18);
  \end{scope}%
}

\providecommand{\HuangHypercube}[3][1]{%
  \begin{scope}[shift={(#2,#3)}, scale=#1]
    \coordinate (z000) at (-1.2,-0.25);
    \coordinate (z001) at (-0.25,0.85);
    \coordinate (z010) at (-0.25,-0.25);
    \coordinate (z011) at (0.65,0.85);
    \coordinate (z100) at (-0.25,-1.2);
    \coordinate (z101) at (0.65,-0.25);
    \coordinate (z110) at (0.65,-1.2);
    \coordinate (z111) at (1.55,-0.25);
    \foreach \a/\b in {z000/z001,z000/z010,z000/z100,z001/z011,z001/z101,z010/z011,z010/z110,z100/z101,z100/z110,z011/z111,z101/z111,z110/z111}
      \draw[line width=1.5pt, huangGray] (\a) -- (\b);
    \foreach \p/\lab/\anch in {z000/000/below left,z001/001/above,z010/010/above left,z011/011/above,z100/100/below,z101/101/above,z110/110/below,z111/111/right}
      \fill[huangRed] (\p) circle (0.08) node[\anch, font=\scriptsize, text=black] {\lab};
  \end{scope}%
}
}

\section{Introduction}

\begin{frame}{Tomography}
    \begin{columns}[T,onlytextwidth]
        \begin{column}{0.56\textwidth}
            \begin{itemize}
                \item<1-> Tomography estimates the full density matrix $\rho \in \mathbb{C}^{d\times d}$
                \item<1-> Number of copies is polynomial in $d$
                \item<2-> With $d=2^n$, this is exponential in the number of qubits.
            \end{itemize}

            \vspace{0.6em}
            \uncover<3->{
                We only need to distinguish high fidelity from low fidelity with one fixed target state, not reconstruct $\rho_{\text{lab}}$
            }
        \end{column}
        \begin{column}{0.38\textwidth}
            \centering
            \begin{tikzpicture}[font=\small, >=Stealth]
                \node[draw, rounded corners=3pt, minimum width=2.25cm, minimum height=0.8cm, fill=blue!8] (copies) at (0,1.2) {$\rho^{\otimes m}$};
                \node[draw, rounded corners=3pt, minimum width=2.25cm, minimum height=0.8cm, fill=orange!12] (tomo) at (0,0) {Tomography};
                \node[draw, rounded corners=3pt, minimum width=2.25cm, minimum height=0.8cm, fill=green!10] (est) at (0,-1.2) {$\widehat{\rho}$};
                \draw[->, thick] (copies) -- (tomo);
                \draw[->, thick] (tomo) -- (est);
                \uncover<2->{\node[align=center, text=red!70!black] at (0,-2.1) {$m$ scales with\\$d=2^n$};}
                \HuangUseBoundingBox{-1.35}{-2.45}{1.35}{1.65}
            \end{tikzpicture}
        \end{column}
    \end{columns}
\end{frame}

\begin{frame}{State Certification}
    \begin{center}
        \resizebox{0.78\textwidth}{!}{\IfFileExists{../diagrams/huang_templates.tex}{\definecolor{huangPurple}{RGB}{113,49,169}
\definecolor{huangBlue}{RGB}{65,109,190}
\definecolor{huangRed}{RGB}{196,0,0}
\definecolor{huangOrange}{RGB}{239,126,45}
\definecolor{huangGold}{RGB}{196,145,15}
\definecolor{huangGreen}{RGB}{91,164,61}
\definecolor{huangGray}{RGB}{120,120,120}
\definecolor{huangPaleBlue}{RGB}{224,234,238}
\definecolor{huangPaleGreen}{RGB}{225,240,216}
\definecolor{huangPeach}{RGB}{246,185,139}
\definecolor{huangLavender}{RGB}{216,205,236}

\tikzset{
  huang arrow/.style={-Stealth, thick, draw=huangGray},
  huang dashed box/.style={draw, dashed, rounded corners=5pt, fill=huangPaleBlue, inner sep=5pt},
  huang bit/.style={draw, dashed, rounded corners=4pt, minimum width=0.42cm, minimum height=0.36cm, inner sep=1pt, font=\scriptsize\bfseries},
  huang gate/.style={draw, rounded corners=2pt, fill=huangPeach, minimum width=0.35cm, minimum height=0.54cm, inner sep=1pt},
  huang meas/.style={draw, rounded corners=2pt, fill=huangLavender, minimum width=0.24cm, minimum height=0.24cm, inner sep=0pt},
  huang label/.style={font=\scriptsize},
  huang title/.style={font=\bfseries}
}

\providecommand{\HuangAtom}[2]{%
  \begin{scope}[shift={(#1,#2)}, scale=0.18]
    \foreach \a in {0,60,...,300} {
      \draw[huangPurple, thick, rotate=\a] (0,0) ellipse (0.36 and 0.9);
    }
    \fill[huangGreen!85!black] (0,0) circle (0.11);
  \end{scope}%
}

\providecommand{\HuangLabPatch}[3][1]{%
  \begin{scope}[shift={(#2,#3)}, scale=#1]
    \node[huang dashed box, minimum width=2.25cm, minimum height=1.55cm] (labpatch) at (0,0) {};
    \foreach \x/\y in {-0.7/0.43,0/0.43,0.7/0.43,-0.7/-0.05,0/-0.05,0.7/-0.05,-0.35/-0.53,0.35/-0.53}
      {\HuangAtom{\x}{\y}}
  \end{scope}%
}

\providecommand{\HuangLabStack}[3][1]{%
  \begin{scope}[shift={(#2,#3)}, scale=#1]
    \foreach \dx in {-0.42,-0.18,0.06} {
      \node[huang dashed box, minimum width=1.75cm, minimum height=1.25cm, fill=huangPaleBlue!85] at (\dx,0) {};
    }
    \foreach \x/\y in {-0.58/0.36,0/0.36,0.58/0.36,-0.58/-0.03,0/-0.03,0.58/-0.03,-0.3/-0.43,0.3/-0.43}
      {\HuangAtom{\x}{\y}}
  \end{scope}%
}

\providecommand{\HuangMonitor}[3][1]{%
  \begin{scope}[shift={(#2,#3)}, scale=#1]
    \fill[black] (-0.72,-0.54) rectangle (0.72,0.55);
    \fill[huangLavender!60] (-0.6,-0.4) rectangle (0.6,0.4);
    \fill[white, opacity=0.72] (-0.55,-0.35) -- (-0.12,0.35) -- (0.6,0.35) -- (0.6,-0.35) -- cycle;
    \node[font=\scriptsize\bfseries, text=huangGray!80] at (0.12,0.02) {Model $\Psi$};
    \fill[black] (0,-0.47) circle (0.06);
    \fill[black] (-0.13,-0.8) rectangle (0.13,-0.54);
    \draw[black, line width=1.5pt] (-0.38,-0.84) -- (0.38,-0.84);
  \end{scope}%
}

\providecommand{\HuangMeasurementStack}[3][1]{%
  \begin{scope}[shift={(#2,#3)}, scale=#1]
    \foreach \i in {0,...,6} {
      \node[huang meas] at (0,0.18*\i) {};
      \draw[huangGold!80!black] (-0.08,0.18*\i-0.02) -- (0.08,0.18*\i+0.07);
    }
  \end{scope}%
}

\providecommand{\HuangCircuitGrid}[4][1]{%
  \begin{scope}[shift={(#2,#3)}, scale=#1]
    \foreach \x in {-0.75,-0.45,-0.15,0.15,0.45,0.75}
      \draw[huangGray] (\x,0.9) -- (\x,-1.05);
    \foreach \x/\y/\w/\h in {
      -0.75/0.55/.34/.45,-0.45/0.18/.24/.42,-0.15/0.55/.24/.52,0.15/0.18/.50/.42,0.45/0.55/.34/.45,0.75/0.18/.24/.42,
      -0.75/-0.25/.24/.46,-0.45/-0.63/.52/.42,-0.15/-0.25/.24/.52,0.15/-0.63/.52/.42,0.45/-0.25/.24/.46,0.75/-0.63/.34/.42}
      \node[huang gate, minimum width=\w cm, minimum height=\h cm] at (\x,\y) {};
    \node[font=\scriptsize] at (0,1.12) {$#4$};
  \end{scope}%
}

\providecommand{\HuangNeuralNet}[3][1]{%
  \begin{scope}[shift={(#2,#3)}, scale=#1]
    \node[draw, dashed, rounded corners=8pt, fill=huangPaleGreen, minimum width=1.65cm, minimum height=2.45cm] at (0,0) {};
    \foreach \x in {-0.45,0,0.45} \fill[huangBlue!38] (\x,0.75) circle (0.08);
    \foreach \x in {-0.58,-0.18,0.22,0.58} \fill[huangBlue!38] (\x,0.15) circle (0.08);
    \foreach \x in {-0.58,-0.18,0.22,0.58} \fill[huangBlue!38] (\x,-0.45) circle (0.08);
    \fill[huangBlue!38] (0,-1.02) circle (0.08);
    \foreach \a in {-0.45,0,0.45} \foreach \b in {-0.58,-0.18,0.22,0.58}
      \draw[huangGray!65, thin] (\a,0.75) -- (\b,0.15);
    \foreach \a in {-0.58,-0.18,0.22,0.58} \foreach \b in {-0.58,-0.18,0.22,0.58}
      \draw[huangGray!65, thin] (\a,0.15) -- (\b,-0.45);
    \foreach \a in {-0.58,-0.18,0.22,0.58}
      \draw[huangGray!65, thin] (\a,-0.45) -- (0,-1.02);
  \end{scope}%
}

\providecommand{\HuangDeviceIcon}[3][1]{%
  \begin{scope}[shift={(#2,#3)}, scale=#1]
    \fill[brown!80!black] (-0.72,0.82) rectangle (0.72,1.12);
    \fill[brown!70!black] (-0.82,0.72) rectangle (0.82,0.78);
    \foreach \x in {-0.48,-0.28,0.28,0.48}
      \draw[line width=1.5pt, huangGold] (\x,0.7) -- (\x,-0.55);
    \foreach \y in {0.42,0.12,-0.18}
      \draw[line width=1.2pt, brown!70!black] (-0.65,\y) -- (0.65,\y);
    \foreach \x in {-0.48,-0.38,0.38,0.48}
      \foreach \y in {0.28,-0.08}
        \draw[line width=1pt, huangGold] (\x,\y) circle (0.055);
    \foreach \x in {-0.52,-0.42,0.42,0.52}
      \foreach \y in {-0.38,-0.48,-0.58}
        \draw[line width=1.5pt, huangGold] (\x-0.08,\y) -- (\x+0.08,\y);
    \fill[huangGold!70] (-0.28,-0.6) rectangle (0.28,-0.85);
    \fill[brown!85!black] (-0.15,-0.85) rectangle (0.15,-1.18);
  \end{scope}%
}

\providecommand{\HuangHypercube}[3][1]{%
  \begin{scope}[shift={(#2,#3)}, scale=#1]
    \coordinate (z000) at (-1.2,-0.25);
    \coordinate (z001) at (-0.25,0.85);
    \coordinate (z010) at (-0.25,-0.25);
    \coordinate (z011) at (0.65,0.85);
    \coordinate (z100) at (-0.25,-1.2);
    \coordinate (z101) at (0.65,-0.25);
    \coordinate (z110) at (0.65,-1.2);
    \coordinate (z111) at (1.55,-0.25);
    \foreach \a/\b in {z000/z001,z000/z010,z000/z100,z001/z011,z001/z101,z010/z011,z010/z110,z100/z101,z100/z110,z011/z111,z101/z111,z110/z111}
      \draw[line width=1.5pt, huangGray] (\a) -- (\b);
    \foreach \p/\lab/\anch in {z000/000/below left,z001/001/above,z010/010/above left,z011/011/above,z100/100/below,z101/101/above,z110/110/below,z111/111/right}
      \fill[huangRed] (\p) circle (0.08) node[\anch, font=\scriptsize, text=black] {\lab};
  \end{scope}%
}
}{\definecolor{huangPurple}{RGB}{113,49,169}
\definecolor{huangBlue}{RGB}{65,109,190}
\definecolor{huangRed}{RGB}{196,0,0}
\definecolor{huangOrange}{RGB}{239,126,45}
\definecolor{huangGold}{RGB}{196,145,15}
\definecolor{huangGreen}{RGB}{91,164,61}
\definecolor{huangGray}{RGB}{120,120,120}
\definecolor{huangPaleBlue}{RGB}{224,234,238}
\definecolor{huangPaleGreen}{RGB}{225,240,216}
\definecolor{huangPeach}{RGB}{246,185,139}
\definecolor{huangLavender}{RGB}{216,205,236}

\tikzset{
  huang arrow/.style={-Stealth, thick, draw=huangGray},
  huang dashed box/.style={draw, dashed, rounded corners=5pt, fill=huangPaleBlue, inner sep=5pt},
  huang bit/.style={draw, dashed, rounded corners=4pt, minimum width=0.42cm, minimum height=0.36cm, inner sep=1pt, font=\scriptsize\bfseries},
  huang gate/.style={draw, rounded corners=2pt, fill=huangPeach, minimum width=0.35cm, minimum height=0.54cm, inner sep=1pt},
  huang meas/.style={draw, rounded corners=2pt, fill=huangLavender, minimum width=0.24cm, minimum height=0.24cm, inner sep=0pt},
  huang label/.style={font=\scriptsize},
  huang title/.style={font=\bfseries}
}

\providecommand{\HuangAtom}[2]{%
  \begin{scope}[shift={(#1,#2)}, scale=0.18]
    \foreach \a in {0,60,...,300} {
      \draw[huangPurple, thick, rotate=\a] (0,0) ellipse (0.36 and 0.9);
    }
    \fill[huangGreen!85!black] (0,0) circle (0.11);
  \end{scope}%
}

\providecommand{\HuangLabPatch}[3][1]{%
  \begin{scope}[shift={(#2,#3)}, scale=#1]
    \node[huang dashed box, minimum width=2.25cm, minimum height=1.55cm] (labpatch) at (0,0) {};
    \foreach \x/\y in {-0.7/0.43,0/0.43,0.7/0.43,-0.7/-0.05,0/-0.05,0.7/-0.05,-0.35/-0.53,0.35/-0.53}
      {\HuangAtom{\x}{\y}}
  \end{scope}%
}

\providecommand{\HuangLabStack}[3][1]{%
  \begin{scope}[shift={(#2,#3)}, scale=#1]
    \foreach \dx in {-0.42,-0.18,0.06} {
      \node[huang dashed box, minimum width=1.75cm, minimum height=1.25cm, fill=huangPaleBlue!85] at (\dx,0) {};
    }
    \foreach \x/\y in {-0.58/0.36,0/0.36,0.58/0.36,-0.58/-0.03,0/-0.03,0.58/-0.03,-0.3/-0.43,0.3/-0.43}
      {\HuangAtom{\x}{\y}}
  \end{scope}%
}

\providecommand{\HuangMonitor}[3][1]{%
  \begin{scope}[shift={(#2,#3)}, scale=#1]
    \fill[black] (-0.72,-0.54) rectangle (0.72,0.55);
    \fill[huangLavender!60] (-0.6,-0.4) rectangle (0.6,0.4);
    \fill[white, opacity=0.72] (-0.55,-0.35) -- (-0.12,0.35) -- (0.6,0.35) -- (0.6,-0.35) -- cycle;
    \node[font=\scriptsize\bfseries, text=huangGray!80] at (0.12,0.02) {Model $\Psi$};
    \fill[black] (0,-0.47) circle (0.06);
    \fill[black] (-0.13,-0.8) rectangle (0.13,-0.54);
    \draw[black, line width=1.5pt] (-0.38,-0.84) -- (0.38,-0.84);
  \end{scope}%
}

\providecommand{\HuangMeasurementStack}[3][1]{%
  \begin{scope}[shift={(#2,#3)}, scale=#1]
    \foreach \i in {0,...,6} {
      \node[huang meas] at (0,0.18*\i) {};
      \draw[huangGold!80!black] (-0.08,0.18*\i-0.02) -- (0.08,0.18*\i+0.07);
    }
  \end{scope}%
}

\providecommand{\HuangCircuitGrid}[4][1]{%
  \begin{scope}[shift={(#2,#3)}, scale=#1]
    \foreach \x in {-0.75,-0.45,-0.15,0.15,0.45,0.75}
      \draw[huangGray] (\x,0.9) -- (\x,-1.05);
    \foreach \x/\y/\w/\h in {
      -0.75/0.55/.34/.45,-0.45/0.18/.24/.42,-0.15/0.55/.24/.52,0.15/0.18/.50/.42,0.45/0.55/.34/.45,0.75/0.18/.24/.42,
      -0.75/-0.25/.24/.46,-0.45/-0.63/.52/.42,-0.15/-0.25/.24/.52,0.15/-0.63/.52/.42,0.45/-0.25/.24/.46,0.75/-0.63/.34/.42}
      \node[huang gate, minimum width=\w cm, minimum height=\h cm] at (\x,\y) {};
    \node[font=\scriptsize] at (0,1.12) {$#4$};
  \end{scope}%
}

\providecommand{\HuangNeuralNet}[3][1]{%
  \begin{scope}[shift={(#2,#3)}, scale=#1]
    \node[draw, dashed, rounded corners=8pt, fill=huangPaleGreen, minimum width=1.65cm, minimum height=2.45cm] at (0,0) {};
    \foreach \x in {-0.45,0,0.45} \fill[huangBlue!38] (\x,0.75) circle (0.08);
    \foreach \x in {-0.58,-0.18,0.22,0.58} \fill[huangBlue!38] (\x,0.15) circle (0.08);
    \foreach \x in {-0.58,-0.18,0.22,0.58} \fill[huangBlue!38] (\x,-0.45) circle (0.08);
    \fill[huangBlue!38] (0,-1.02) circle (0.08);
    \foreach \a in {-0.45,0,0.45} \foreach \b in {-0.58,-0.18,0.22,0.58}
      \draw[huangGray!65, thin] (\a,0.75) -- (\b,0.15);
    \foreach \a in {-0.58,-0.18,0.22,0.58} \foreach \b in {-0.58,-0.18,0.22,0.58}
      \draw[huangGray!65, thin] (\a,0.15) -- (\b,-0.45);
    \foreach \a in {-0.58,-0.18,0.22,0.58}
      \draw[huangGray!65, thin] (\a,-0.45) -- (0,-1.02);
  \end{scope}%
}

\providecommand{\HuangDeviceIcon}[3][1]{%
  \begin{scope}[shift={(#2,#3)}, scale=#1]
    \fill[brown!80!black] (-0.72,0.82) rectangle (0.72,1.12);
    \fill[brown!70!black] (-0.82,0.72) rectangle (0.82,0.78);
    \foreach \x in {-0.48,-0.28,0.28,0.48}
      \draw[line width=1.5pt, huangGold] (\x,0.7) -- (\x,-0.55);
    \foreach \y in {0.42,0.12,-0.18}
      \draw[line width=1.2pt, brown!70!black] (-0.65,\y) -- (0.65,\y);
    \foreach \x in {-0.48,-0.38,0.38,0.48}
      \foreach \y in {0.28,-0.08}
        \draw[line width=1pt, huangGold] (\x,\y) circle (0.055);
    \foreach \x in {-0.52,-0.42,0.42,0.52}
      \foreach \y in {-0.38,-0.48,-0.58}
        \draw[line width=1.5pt, huangGold] (\x-0.08,\y) -- (\x+0.08,\y);
    \fill[huangGold!70] (-0.28,-0.6) rectangle (0.28,-0.85);
    \fill[brown!85!black] (-0.15,-0.85) rectangle (0.15,-1.18);
  \end{scope}%
}

\providecommand{\HuangHypercube}[3][1]{%
  \begin{scope}[shift={(#2,#3)}, scale=#1]
    \coordinate (z000) at (-1.2,-0.25);
    \coordinate (z001) at (-0.25,0.85);
    \coordinate (z010) at (-0.25,-0.25);
    \coordinate (z011) at (0.65,0.85);
    \coordinate (z100) at (-0.25,-1.2);
    \coordinate (z101) at (0.65,-0.25);
    \coordinate (z110) at (0.65,-1.2);
    \coordinate (z111) at (1.55,-0.25);
    \foreach \a/\b in {z000/z001,z000/z010,z000/z100,z001/z011,z001/z101,z010/z011,z010/z110,z100/z101,z100/z110,z011/z111,z101/z111,z110/z111}
      \draw[line width=1.5pt, huangGray] (\a) -- (\b);
    \foreach \p/\lab/\anch in {z000/000/below left,z001/001/above,z010/010/above left,z011/011/above,z100/100/below,z101/101/above,z110/110/below,z111/111/right}
      \fill[huangRed] (\p) circle (0.08) node[\anch, font=\scriptsize, text=black] {\lab};
  \end{scope}%
}
}
\begin{tikzpicture}[>=Stealth, font=\footnotesize]
\node[font=\bfseries] at (-3.2,1.8) {Lab state $\rho$};
\HuangLabStack[1.25]{-3.25}{0.65}
\node[text=huangGray, font=\bfseries] at (-3.2,-0.55) {measurement data};

\node[font=\bfseries] at (3.0,1.8) {Target state};
\node at (3.0,1.45) {$|\psi\rangle=\sum_{x\in\{0,1\}^n}\psi(x)|x\rangle$};
\HuangMonitor[1.0]{3.0}{0.45}
\draw[huang arrow] (1.0,0.55) -- (2.05,0.55) node[midway, above] {$x\in\{0,1\}^n$};
\draw[huang arrow] (3.95,0.55) -- (4.85,0.55) node[midway, above, text=huangRed] {$\psi(x)$};
\node[text=huangGray, font=\bfseries] at (3.0,-0.65) {query access to $|\psi\rangle$};

\draw[huang arrow] (-3.2,-0.85) -- (-2.25,-1.55);
\draw[huang arrow] (3.0,-0.9) -- (2.1,-1.55);
\node[font=\large] at (0,-1.95) {$\langle\psi|\rho|\psi\rangle$ is close to $1$ or far from $1$?};

\HuangUseBoundingBox{-5.4}{-2.25}{5.4}{2.1}

\end{tikzpicture}
}
    \end{center}

    \vspace{-0.3em}
    \begin{block}{Goal}
        Given a classical description of a target pure state $\ket{\psi_{\text{target}}}$ and one copy of an unknown lab state $\rho_{\text{lab}}$, decide whether
        \[
            \mathrm{Fid}=\bra{\psi_{\text{target}}}\rho_{\text{lab}}\ket{\psi_{\text{target}}}
        \]
        is large, using only simple measurements on $\rho_{\text{lab}}$.
    \end{block}
\end{frame}

\begin{frame}{Naive Certification}
    \begin{columns}[T,onlytextwidth]
        \begin{column}{0.5\textwidth}
            It's trivial! Just use
            \[
                \left\{\ket{\psi_{\text{target}}}\!\bra{\psi_{\text{target}}},
                I-\ket{\psi_{\text{target}}}\!\bra{\psi_{\text{target}}}\right\}.
            \]

            \pause
            \vspace{0.5em}
            Equivalently, one may use a SWAP test against a freshly prepared copy of $\ket{\psi_{\text{target}}}$.
    \pause

    \vspace{0.3em}
      Circular: We need $\ket{\psi_{\text{target}}}$ to test $\ket{\psi_{\text{target}}}$
        \end{column}
        \begin{column}{0.46\textwidth}
            \centering
            \resizebox{\linewidth}{!}{\IfFileExists{../diagrams/huang_templates.tex}{\definecolor{huangPurple}{RGB}{113,49,169}
\definecolor{huangBlue}{RGB}{65,109,190}
\definecolor{huangRed}{RGB}{196,0,0}
\definecolor{huangOrange}{RGB}{239,126,45}
\definecolor{huangGold}{RGB}{196,145,15}
\definecolor{huangGreen}{RGB}{91,164,61}
\definecolor{huangGray}{RGB}{120,120,120}
\definecolor{huangPaleBlue}{RGB}{224,234,238}
\definecolor{huangPaleGreen}{RGB}{225,240,216}
\definecolor{huangPeach}{RGB}{246,185,139}
\definecolor{huangLavender}{RGB}{216,205,236}

\tikzset{
  huang arrow/.style={-Stealth, thick, draw=huangGray},
  huang dashed box/.style={draw, dashed, rounded corners=5pt, fill=huangPaleBlue, inner sep=5pt},
  huang bit/.style={draw, dashed, rounded corners=4pt, minimum width=0.42cm, minimum height=0.36cm, inner sep=1pt, font=\scriptsize\bfseries},
  huang gate/.style={draw, rounded corners=2pt, fill=huangPeach, minimum width=0.35cm, minimum height=0.54cm, inner sep=1pt},
  huang meas/.style={draw, rounded corners=2pt, fill=huangLavender, minimum width=0.24cm, minimum height=0.24cm, inner sep=0pt},
  huang label/.style={font=\scriptsize},
  huang title/.style={font=\bfseries}
}

\providecommand{\HuangAtom}[2]{%
  \begin{scope}[shift={(#1,#2)}, scale=0.18]
    \foreach \a in {0,60,...,300} {
      \draw[huangPurple, thick, rotate=\a] (0,0) ellipse (0.36 and 0.9);
    }
    \fill[huangGreen!85!black] (0,0) circle (0.11);
  \end{scope}%
}

\providecommand{\HuangLabPatch}[3][1]{%
  \begin{scope}[shift={(#2,#3)}, scale=#1]
    \node[huang dashed box, minimum width=2.25cm, minimum height=1.55cm] (labpatch) at (0,0) {};
    \foreach \x/\y in {-0.7/0.43,0/0.43,0.7/0.43,-0.7/-0.05,0/-0.05,0.7/-0.05,-0.35/-0.53,0.35/-0.53}
      {\HuangAtom{\x}{\y}}
  \end{scope}%
}

\providecommand{\HuangLabStack}[3][1]{%
  \begin{scope}[shift={(#2,#3)}, scale=#1]
    \foreach \dx in {-0.42,-0.18,0.06} {
      \node[huang dashed box, minimum width=1.75cm, minimum height=1.25cm, fill=huangPaleBlue!85] at (\dx,0) {};
    }
    \foreach \x/\y in {-0.58/0.36,0/0.36,0.58/0.36,-0.58/-0.03,0/-0.03,0.58/-0.03,-0.3/-0.43,0.3/-0.43}
      {\HuangAtom{\x}{\y}}
  \end{scope}%
}

\providecommand{\HuangMonitor}[3][1]{%
  \begin{scope}[shift={(#2,#3)}, scale=#1]
    \fill[black] (-0.72,-0.54) rectangle (0.72,0.55);
    \fill[huangLavender!60] (-0.6,-0.4) rectangle (0.6,0.4);
    \fill[white, opacity=0.72] (-0.55,-0.35) -- (-0.12,0.35) -- (0.6,0.35) -- (0.6,-0.35) -- cycle;
    \node[font=\scriptsize\bfseries, text=huangGray!80] at (0.12,0.02) {Model $\Psi$};
    \fill[black] (0,-0.47) circle (0.06);
    \fill[black] (-0.13,-0.8) rectangle (0.13,-0.54);
    \draw[black, line width=1.5pt] (-0.38,-0.84) -- (0.38,-0.84);
  \end{scope}%
}

\providecommand{\HuangMeasurementStack}[3][1]{%
  \begin{scope}[shift={(#2,#3)}, scale=#1]
    \foreach \i in {0,...,6} {
      \node[huang meas] at (0,0.18*\i) {};
      \draw[huangGold!80!black] (-0.08,0.18*\i-0.02) -- (0.08,0.18*\i+0.07);
    }
  \end{scope}%
}

\providecommand{\HuangCircuitGrid}[4][1]{%
  \begin{scope}[shift={(#2,#3)}, scale=#1]
    \foreach \x in {-0.75,-0.45,-0.15,0.15,0.45,0.75}
      \draw[huangGray] (\x,0.9) -- (\x,-1.05);
    \foreach \x/\y/\w/\h in {
      -0.75/0.55/.34/.45,-0.45/0.18/.24/.42,-0.15/0.55/.24/.52,0.15/0.18/.50/.42,0.45/0.55/.34/.45,0.75/0.18/.24/.42,
      -0.75/-0.25/.24/.46,-0.45/-0.63/.52/.42,-0.15/-0.25/.24/.52,0.15/-0.63/.52/.42,0.45/-0.25/.24/.46,0.75/-0.63/.34/.42}
      \node[huang gate, minimum width=\w cm, minimum height=\h cm] at (\x,\y) {};
    \node[font=\scriptsize] at (0,1.12) {$#4$};
  \end{scope}%
}

\providecommand{\HuangNeuralNet}[3][1]{%
  \begin{scope}[shift={(#2,#3)}, scale=#1]
    \node[draw, dashed, rounded corners=8pt, fill=huangPaleGreen, minimum width=1.65cm, minimum height=2.45cm] at (0,0) {};
    \foreach \x in {-0.45,0,0.45} \fill[huangBlue!38] (\x,0.75) circle (0.08);
    \foreach \x in {-0.58,-0.18,0.22,0.58} \fill[huangBlue!38] (\x,0.15) circle (0.08);
    \foreach \x in {-0.58,-0.18,0.22,0.58} \fill[huangBlue!38] (\x,-0.45) circle (0.08);
    \fill[huangBlue!38] (0,-1.02) circle (0.08);
    \foreach \a in {-0.45,0,0.45} \foreach \b in {-0.58,-0.18,0.22,0.58}
      \draw[huangGray!65, thin] (\a,0.75) -- (\b,0.15);
    \foreach \a in {-0.58,-0.18,0.22,0.58} \foreach \b in {-0.58,-0.18,0.22,0.58}
      \draw[huangGray!65, thin] (\a,0.15) -- (\b,-0.45);
    \foreach \a in {-0.58,-0.18,0.22,0.58}
      \draw[huangGray!65, thin] (\a,-0.45) -- (0,-1.02);
  \end{scope}%
}

\providecommand{\HuangDeviceIcon}[3][1]{%
  \begin{scope}[shift={(#2,#3)}, scale=#1]
    \fill[brown!80!black] (-0.72,0.82) rectangle (0.72,1.12);
    \fill[brown!70!black] (-0.82,0.72) rectangle (0.82,0.78);
    \foreach \x in {-0.48,-0.28,0.28,0.48}
      \draw[line width=1.5pt, huangGold] (\x,0.7) -- (\x,-0.55);
    \foreach \y in {0.42,0.12,-0.18}
      \draw[line width=1.2pt, brown!70!black] (-0.65,\y) -- (0.65,\y);
    \foreach \x in {-0.48,-0.38,0.38,0.48}
      \foreach \y in {0.28,-0.08}
        \draw[line width=1pt, huangGold] (\x,\y) circle (0.055);
    \foreach \x in {-0.52,-0.42,0.42,0.52}
      \foreach \y in {-0.38,-0.48,-0.58}
        \draw[line width=1.5pt, huangGold] (\x-0.08,\y) -- (\x+0.08,\y);
    \fill[huangGold!70] (-0.28,-0.6) rectangle (0.28,-0.85);
    \fill[brown!85!black] (-0.15,-0.85) rectangle (0.15,-1.18);
  \end{scope}%
}

\providecommand{\HuangHypercube}[3][1]{%
  \begin{scope}[shift={(#2,#3)}, scale=#1]
    \coordinate (z000) at (-1.2,-0.25);
    \coordinate (z001) at (-0.25,0.85);
    \coordinate (z010) at (-0.25,-0.25);
    \coordinate (z011) at (0.65,0.85);
    \coordinate (z100) at (-0.25,-1.2);
    \coordinate (z101) at (0.65,-0.25);
    \coordinate (z110) at (0.65,-1.2);
    \coordinate (z111) at (1.55,-0.25);
    \foreach \a/\b in {z000/z001,z000/z010,z000/z100,z001/z011,z001/z101,z010/z011,z010/z110,z100/z101,z100/z110,z011/z111,z101/z111,z110/z111}
      \draw[line width=1.5pt, huangGray] (\a) -- (\b);
    \foreach \p/\lab/\anch in {z000/000/below left,z001/001/above,z010/010/above left,z011/011/above,z100/100/below,z101/101/above,z110/110/below,z111/111/right}
      \fill[huangRed] (\p) circle (0.08) node[\anch, font=\scriptsize, text=black] {\lab};
  \end{scope}%
}
}{\definecolor{huangPurple}{RGB}{113,49,169}
\definecolor{huangBlue}{RGB}{65,109,190}
\definecolor{huangRed}{RGB}{196,0,0}
\definecolor{huangOrange}{RGB}{239,126,45}
\definecolor{huangGold}{RGB}{196,145,15}
\definecolor{huangGreen}{RGB}{91,164,61}
\definecolor{huangGray}{RGB}{120,120,120}
\definecolor{huangPaleBlue}{RGB}{224,234,238}
\definecolor{huangPaleGreen}{RGB}{225,240,216}
\definecolor{huangPeach}{RGB}{246,185,139}
\definecolor{huangLavender}{RGB}{216,205,236}

\tikzset{
  huang arrow/.style={-Stealth, thick, draw=huangGray},
  huang dashed box/.style={draw, dashed, rounded corners=5pt, fill=huangPaleBlue, inner sep=5pt},
  huang bit/.style={draw, dashed, rounded corners=4pt, minimum width=0.42cm, minimum height=0.36cm, inner sep=1pt, font=\scriptsize\bfseries},
  huang gate/.style={draw, rounded corners=2pt, fill=huangPeach, minimum width=0.35cm, minimum height=0.54cm, inner sep=1pt},
  huang meas/.style={draw, rounded corners=2pt, fill=huangLavender, minimum width=0.24cm, minimum height=0.24cm, inner sep=0pt},
  huang label/.style={font=\scriptsize},
  huang title/.style={font=\bfseries}
}

\providecommand{\HuangAtom}[2]{%
  \begin{scope}[shift={(#1,#2)}, scale=0.18]
    \foreach \a in {0,60,...,300} {
      \draw[huangPurple, thick, rotate=\a] (0,0) ellipse (0.36 and 0.9);
    }
    \fill[huangGreen!85!black] (0,0) circle (0.11);
  \end{scope}%
}

\providecommand{\HuangLabPatch}[3][1]{%
  \begin{scope}[shift={(#2,#3)}, scale=#1]
    \node[huang dashed box, minimum width=2.25cm, minimum height=1.55cm] (labpatch) at (0,0) {};
    \foreach \x/\y in {-0.7/0.43,0/0.43,0.7/0.43,-0.7/-0.05,0/-0.05,0.7/-0.05,-0.35/-0.53,0.35/-0.53}
      {\HuangAtom{\x}{\y}}
  \end{scope}%
}

\providecommand{\HuangLabStack}[3][1]{%
  \begin{scope}[shift={(#2,#3)}, scale=#1]
    \foreach \dx in {-0.42,-0.18,0.06} {
      \node[huang dashed box, minimum width=1.75cm, minimum height=1.25cm, fill=huangPaleBlue!85] at (\dx,0) {};
    }
    \foreach \x/\y in {-0.58/0.36,0/0.36,0.58/0.36,-0.58/-0.03,0/-0.03,0.58/-0.03,-0.3/-0.43,0.3/-0.43}
      {\HuangAtom{\x}{\y}}
  \end{scope}%
}

\providecommand{\HuangMonitor}[3][1]{%
  \begin{scope}[shift={(#2,#3)}, scale=#1]
    \fill[black] (-0.72,-0.54) rectangle (0.72,0.55);
    \fill[huangLavender!60] (-0.6,-0.4) rectangle (0.6,0.4);
    \fill[white, opacity=0.72] (-0.55,-0.35) -- (-0.12,0.35) -- (0.6,0.35) -- (0.6,-0.35) -- cycle;
    \node[font=\scriptsize\bfseries, text=huangGray!80] at (0.12,0.02) {Model $\Psi$};
    \fill[black] (0,-0.47) circle (0.06);
    \fill[black] (-0.13,-0.8) rectangle (0.13,-0.54);
    \draw[black, line width=1.5pt] (-0.38,-0.84) -- (0.38,-0.84);
  \end{scope}%
}

\providecommand{\HuangMeasurementStack}[3][1]{%
  \begin{scope}[shift={(#2,#3)}, scale=#1]
    \foreach \i in {0,...,6} {
      \node[huang meas] at (0,0.18*\i) {};
      \draw[huangGold!80!black] (-0.08,0.18*\i-0.02) -- (0.08,0.18*\i+0.07);
    }
  \end{scope}%
}

\providecommand{\HuangCircuitGrid}[4][1]{%
  \begin{scope}[shift={(#2,#3)}, scale=#1]
    \foreach \x in {-0.75,-0.45,-0.15,0.15,0.45,0.75}
      \draw[huangGray] (\x,0.9) -- (\x,-1.05);
    \foreach \x/\y/\w/\h in {
      -0.75/0.55/.34/.45,-0.45/0.18/.24/.42,-0.15/0.55/.24/.52,0.15/0.18/.50/.42,0.45/0.55/.34/.45,0.75/0.18/.24/.42,
      -0.75/-0.25/.24/.46,-0.45/-0.63/.52/.42,-0.15/-0.25/.24/.52,0.15/-0.63/.52/.42,0.45/-0.25/.24/.46,0.75/-0.63/.34/.42}
      \node[huang gate, minimum width=\w cm, minimum height=\h cm] at (\x,\y) {};
    \node[font=\scriptsize] at (0,1.12) {$#4$};
  \end{scope}%
}

\providecommand{\HuangNeuralNet}[3][1]{%
  \begin{scope}[shift={(#2,#3)}, scale=#1]
    \node[draw, dashed, rounded corners=8pt, fill=huangPaleGreen, minimum width=1.65cm, minimum height=2.45cm] at (0,0) {};
    \foreach \x in {-0.45,0,0.45} \fill[huangBlue!38] (\x,0.75) circle (0.08);
    \foreach \x in {-0.58,-0.18,0.22,0.58} \fill[huangBlue!38] (\x,0.15) circle (0.08);
    \foreach \x in {-0.58,-0.18,0.22,0.58} \fill[huangBlue!38] (\x,-0.45) circle (0.08);
    \fill[huangBlue!38] (0,-1.02) circle (0.08);
    \foreach \a in {-0.45,0,0.45} \foreach \b in {-0.58,-0.18,0.22,0.58}
      \draw[huangGray!65, thin] (\a,0.75) -- (\b,0.15);
    \foreach \a in {-0.58,-0.18,0.22,0.58} \foreach \b in {-0.58,-0.18,0.22,0.58}
      \draw[huangGray!65, thin] (\a,0.15) -- (\b,-0.45);
    \foreach \a in {-0.58,-0.18,0.22,0.58}
      \draw[huangGray!65, thin] (\a,-0.45) -- (0,-1.02);
  \end{scope}%
}

\providecommand{\HuangDeviceIcon}[3][1]{%
  \begin{scope}[shift={(#2,#3)}, scale=#1]
    \fill[brown!80!black] (-0.72,0.82) rectangle (0.72,1.12);
    \fill[brown!70!black] (-0.82,0.72) rectangle (0.82,0.78);
    \foreach \x in {-0.48,-0.28,0.28,0.48}
      \draw[line width=1.5pt, huangGold] (\x,0.7) -- (\x,-0.55);
    \foreach \y in {0.42,0.12,-0.18}
      \draw[line width=1.2pt, brown!70!black] (-0.65,\y) -- (0.65,\y);
    \foreach \x in {-0.48,-0.38,0.38,0.48}
      \foreach \y in {0.28,-0.08}
        \draw[line width=1pt, huangGold] (\x,\y) circle (0.055);
    \foreach \x in {-0.52,-0.42,0.42,0.52}
      \foreach \y in {-0.38,-0.48,-0.58}
        \draw[line width=1.5pt, huangGold] (\x-0.08,\y) -- (\x+0.08,\y);
    \fill[huangGold!70] (-0.28,-0.6) rectangle (0.28,-0.85);
    \fill[brown!85!black] (-0.15,-0.85) rectangle (0.15,-1.18);
  \end{scope}%
}

\providecommand{\HuangHypercube}[3][1]{%
  \begin{scope}[shift={(#2,#3)}, scale=#1]
    \coordinate (z000) at (-1.2,-0.25);
    \coordinate (z001) at (-0.25,0.85);
    \coordinate (z010) at (-0.25,-0.25);
    \coordinate (z011) at (0.65,0.85);
    \coordinate (z100) at (-0.25,-1.2);
    \coordinate (z101) at (0.65,-0.25);
    \coordinate (z110) at (0.65,-1.2);
    \coordinate (z111) at (1.55,-0.25);
    \foreach \a/\b in {z000/z001,z000/z010,z000/z100,z001/z011,z001/z101,z010/z011,z010/z110,z100/z101,z100/z110,z011/z111,z101/z111,z110/z111}
      \draw[line width=1.5pt, huangGray] (\a) -- (\b);
    \foreach \p/\lab/\anch in {z000/000/below left,z001/001/above,z010/010/above left,z011/011/above,z100/100/below,z101/101/above,z110/110/below,z111/111/right}
      \fill[huangRed] (\p) circle (0.08) node[\anch, font=\scriptsize, text=black] {\lab};
  \end{scope}%
}
}
\begin{tikzpicture}[>=Stealth, font=\footnotesize]
\node[font=\bfseries] at (-2.8,1.2) {Lab state $\rho$};
\HuangLabPatch[0.85]{-2.8}{0.35}
\foreach \y in {-0.45,-0.15,0.15,0.45} \draw[huangGray] (-1.65,\y) -- (-1.25,\y);
\node[draw=huangBlue, dashed, rounded corners=4pt, fill=brown!25,
  minimum width=1.95cm, minimum height=1.45cm, font=\large] at (0,0) {$U^\dagger$};
\foreach \y in {-0.54,-0.36,-0.18,0,0.18,0.36,0.54}
  \node[huang meas] at (1.25,\y) {};
\draw[huang arrow] (1.45,0) -- (2.3,0) node[right] {$|0^n\rangle$?};
\node[align=center] at (0,-1.35)
  {Complex inverse circuit or SWAP test\\for general target $|\psi\rangle=U|0^n\rangle$};

\HuangUseBoundingBox{-3.95}{-1.75}{3.95}{1.55}

\end{tikzpicture}
}
        \end{column}
    \end{columns}

\end{frame}

\begin{frame}{Classical Shadows}
    \begin{columns}[T,onlytextwidth]
        \begin{column}{0.48\textwidth}
            \begin{center}
                \makebox[0pt][c]{\hspace*{-2.5cm}\resizebox{0.95\linewidth}{!}{\IfFileExists{../diagrams/huang_templates.tex}{\definecolor{huangPurple}{RGB}{113,49,169}
\definecolor{huangBlue}{RGB}{65,109,190}
\definecolor{huangRed}{RGB}{196,0,0}
\definecolor{huangOrange}{RGB}{239,126,45}
\definecolor{huangGold}{RGB}{196,145,15}
\definecolor{huangGreen}{RGB}{91,164,61}
\definecolor{huangGray}{RGB}{120,120,120}
\definecolor{huangPaleBlue}{RGB}{224,234,238}
\definecolor{huangPaleGreen}{RGB}{225,240,216}
\definecolor{huangPeach}{RGB}{246,185,139}
\definecolor{huangLavender}{RGB}{216,205,236}

\tikzset{
  huang arrow/.style={-Stealth, thick, draw=huangGray},
  huang dashed box/.style={draw, dashed, rounded corners=5pt, fill=huangPaleBlue, inner sep=5pt},
  huang bit/.style={draw, dashed, rounded corners=4pt, minimum width=0.42cm, minimum height=0.36cm, inner sep=1pt, font=\scriptsize\bfseries},
  huang gate/.style={draw, rounded corners=2pt, fill=huangPeach, minimum width=0.35cm, minimum height=0.54cm, inner sep=1pt},
  huang meas/.style={draw, rounded corners=2pt, fill=huangLavender, minimum width=0.24cm, minimum height=0.24cm, inner sep=0pt},
  huang label/.style={font=\scriptsize},
  huang title/.style={font=\bfseries}
}

\providecommand{\HuangAtom}[2]{%
  \begin{scope}[shift={(#1,#2)}, scale=0.18]
    \foreach \a in {0,60,...,300} {
      \draw[huangPurple, thick, rotate=\a] (0,0) ellipse (0.36 and 0.9);
    }
    \fill[huangGreen!85!black] (0,0) circle (0.11);
  \end{scope}%
}

\providecommand{\HuangLabPatch}[3][1]{%
  \begin{scope}[shift={(#2,#3)}, scale=#1]
    \node[huang dashed box, minimum width=2.25cm, minimum height=1.55cm] (labpatch) at (0,0) {};
    \foreach \x/\y in {-0.7/0.43,0/0.43,0.7/0.43,-0.7/-0.05,0/-0.05,0.7/-0.05,-0.35/-0.53,0.35/-0.53}
      {\HuangAtom{\x}{\y}}
  \end{scope}%
}

\providecommand{\HuangLabStack}[3][1]{%
  \begin{scope}[shift={(#2,#3)}, scale=#1]
    \foreach \dx in {-0.42,-0.18,0.06} {
      \node[huang dashed box, minimum width=1.75cm, minimum height=1.25cm, fill=huangPaleBlue!85] at (\dx,0) {};
    }
    \foreach \x/\y in {-0.58/0.36,0/0.36,0.58/0.36,-0.58/-0.03,0/-0.03,0.58/-0.03,-0.3/-0.43,0.3/-0.43}
      {\HuangAtom{\x}{\y}}
  \end{scope}%
}

\providecommand{\HuangMonitor}[3][1]{%
  \begin{scope}[shift={(#2,#3)}, scale=#1]
    \fill[black] (-0.72,-0.54) rectangle (0.72,0.55);
    \fill[huangLavender!60] (-0.6,-0.4) rectangle (0.6,0.4);
    \fill[white, opacity=0.72] (-0.55,-0.35) -- (-0.12,0.35) -- (0.6,0.35) -- (0.6,-0.35) -- cycle;
    \node[font=\scriptsize\bfseries, text=huangGray!80] at (0.12,0.02) {Model $\Psi$};
    \fill[black] (0,-0.47) circle (0.06);
    \fill[black] (-0.13,-0.8) rectangle (0.13,-0.54);
    \draw[black, line width=1.5pt] (-0.38,-0.84) -- (0.38,-0.84);
  \end{scope}%
}

\providecommand{\HuangMeasurementStack}[3][1]{%
  \begin{scope}[shift={(#2,#3)}, scale=#1]
    \foreach \i in {0,...,6} {
      \node[huang meas] at (0,0.18*\i) {};
      \draw[huangGold!80!black] (-0.08,0.18*\i-0.02) -- (0.08,0.18*\i+0.07);
    }
  \end{scope}%
}

\providecommand{\HuangCircuitGrid}[4][1]{%
  \begin{scope}[shift={(#2,#3)}, scale=#1]
    \foreach \x in {-0.75,-0.45,-0.15,0.15,0.45,0.75}
      \draw[huangGray] (\x,0.9) -- (\x,-1.05);
    \foreach \x/\y/\w/\h in {
      -0.75/0.55/.34/.45,-0.45/0.18/.24/.42,-0.15/0.55/.24/.52,0.15/0.18/.50/.42,0.45/0.55/.34/.45,0.75/0.18/.24/.42,
      -0.75/-0.25/.24/.46,-0.45/-0.63/.52/.42,-0.15/-0.25/.24/.52,0.15/-0.63/.52/.42,0.45/-0.25/.24/.46,0.75/-0.63/.34/.42}
      \node[huang gate, minimum width=\w cm, minimum height=\h cm] at (\x,\y) {};
    \node[font=\scriptsize] at (0,1.12) {$#4$};
  \end{scope}%
}

\providecommand{\HuangNeuralNet}[3][1]{%
  \begin{scope}[shift={(#2,#3)}, scale=#1]
    \node[draw, dashed, rounded corners=8pt, fill=huangPaleGreen, minimum width=1.65cm, minimum height=2.45cm] at (0,0) {};
    \foreach \x in {-0.45,0,0.45} \fill[huangBlue!38] (\x,0.75) circle (0.08);
    \foreach \x in {-0.58,-0.18,0.22,0.58} \fill[huangBlue!38] (\x,0.15) circle (0.08);
    \foreach \x in {-0.58,-0.18,0.22,0.58} \fill[huangBlue!38] (\x,-0.45) circle (0.08);
    \fill[huangBlue!38] (0,-1.02) circle (0.08);
    \foreach \a in {-0.45,0,0.45} \foreach \b in {-0.58,-0.18,0.22,0.58}
      \draw[huangGray!65, thin] (\a,0.75) -- (\b,0.15);
    \foreach \a in {-0.58,-0.18,0.22,0.58} \foreach \b in {-0.58,-0.18,0.22,0.58}
      \draw[huangGray!65, thin] (\a,0.15) -- (\b,-0.45);
    \foreach \a in {-0.58,-0.18,0.22,0.58}
      \draw[huangGray!65, thin] (\a,-0.45) -- (0,-1.02);
  \end{scope}%
}

\providecommand{\HuangDeviceIcon}[3][1]{%
  \begin{scope}[shift={(#2,#3)}, scale=#1]
    \fill[brown!80!black] (-0.72,0.82) rectangle (0.72,1.12);
    \fill[brown!70!black] (-0.82,0.72) rectangle (0.82,0.78);
    \foreach \x in {-0.48,-0.28,0.28,0.48}
      \draw[line width=1.5pt, huangGold] (\x,0.7) -- (\x,-0.55);
    \foreach \y in {0.42,0.12,-0.18}
      \draw[line width=1.2pt, brown!70!black] (-0.65,\y) -- (0.65,\y);
    \foreach \x in {-0.48,-0.38,0.38,0.48}
      \foreach \y in {0.28,-0.08}
        \draw[line width=1pt, huangGold] (\x,\y) circle (0.055);
    \foreach \x in {-0.52,-0.42,0.42,0.52}
      \foreach \y in {-0.38,-0.48,-0.58}
        \draw[line width=1.5pt, huangGold] (\x-0.08,\y) -- (\x+0.08,\y);
    \fill[huangGold!70] (-0.28,-0.6) rectangle (0.28,-0.85);
    \fill[brown!85!black] (-0.15,-0.85) rectangle (0.15,-1.18);
  \end{scope}%
}

\providecommand{\HuangHypercube}[3][1]{%
  \begin{scope}[shift={(#2,#3)}, scale=#1]
    \coordinate (z000) at (-1.2,-0.25);
    \coordinate (z001) at (-0.25,0.85);
    \coordinate (z010) at (-0.25,-0.25);
    \coordinate (z011) at (0.65,0.85);
    \coordinate (z100) at (-0.25,-1.2);
    \coordinate (z101) at (0.65,-0.25);
    \coordinate (z110) at (0.65,-1.2);
    \coordinate (z111) at (1.55,-0.25);
    \foreach \a/\b in {z000/z001,z000/z010,z000/z100,z001/z011,z001/z101,z010/z011,z010/z110,z100/z101,z100/z110,z011/z111,z101/z111,z110/z111}
      \draw[line width=1.5pt, huangGray] (\a) -- (\b);
    \foreach \p/\lab/\anch in {z000/000/below left,z001/001/above,z010/010/above left,z011/011/above,z100/100/below,z101/101/above,z110/110/below,z111/111/right}
      \fill[huangRed] (\p) circle (0.08) node[\anch, font=\scriptsize, text=black] {\lab};
  \end{scope}%
}
}{\definecolor{huangPurple}{RGB}{113,49,169}
\definecolor{huangBlue}{RGB}{65,109,190}
\definecolor{huangRed}{RGB}{196,0,0}
\definecolor{huangOrange}{RGB}{239,126,45}
\definecolor{huangGold}{RGB}{196,145,15}
\definecolor{huangGreen}{RGB}{91,164,61}
\definecolor{huangGray}{RGB}{120,120,120}
\definecolor{huangPaleBlue}{RGB}{224,234,238}
\definecolor{huangPaleGreen}{RGB}{225,240,216}
\definecolor{huangPeach}{RGB}{246,185,139}
\definecolor{huangLavender}{RGB}{216,205,236}

\tikzset{
  huang arrow/.style={-Stealth, thick, draw=huangGray},
  huang dashed box/.style={draw, dashed, rounded corners=5pt, fill=huangPaleBlue, inner sep=5pt},
  huang bit/.style={draw, dashed, rounded corners=4pt, minimum width=0.42cm, minimum height=0.36cm, inner sep=1pt, font=\scriptsize\bfseries},
  huang gate/.style={draw, rounded corners=2pt, fill=huangPeach, minimum width=0.35cm, minimum height=0.54cm, inner sep=1pt},
  huang meas/.style={draw, rounded corners=2pt, fill=huangLavender, minimum width=0.24cm, minimum height=0.24cm, inner sep=0pt},
  huang label/.style={font=\scriptsize},
  huang title/.style={font=\bfseries}
}

\providecommand{\HuangAtom}[2]{%
  \begin{scope}[shift={(#1,#2)}, scale=0.18]
    \foreach \a in {0,60,...,300} {
      \draw[huangPurple, thick, rotate=\a] (0,0) ellipse (0.36 and 0.9);
    }
    \fill[huangGreen!85!black] (0,0) circle (0.11);
  \end{scope}%
}

\providecommand{\HuangLabPatch}[3][1]{%
  \begin{scope}[shift={(#2,#3)}, scale=#1]
    \node[huang dashed box, minimum width=2.25cm, minimum height=1.55cm] (labpatch) at (0,0) {};
    \foreach \x/\y in {-0.7/0.43,0/0.43,0.7/0.43,-0.7/-0.05,0/-0.05,0.7/-0.05,-0.35/-0.53,0.35/-0.53}
      {\HuangAtom{\x}{\y}}
  \end{scope}%
}

\providecommand{\HuangLabStack}[3][1]{%
  \begin{scope}[shift={(#2,#3)}, scale=#1]
    \foreach \dx in {-0.42,-0.18,0.06} {
      \node[huang dashed box, minimum width=1.75cm, minimum height=1.25cm, fill=huangPaleBlue!85] at (\dx,0) {};
    }
    \foreach \x/\y in {-0.58/0.36,0/0.36,0.58/0.36,-0.58/-0.03,0/-0.03,0.58/-0.03,-0.3/-0.43,0.3/-0.43}
      {\HuangAtom{\x}{\y}}
  \end{scope}%
}

\providecommand{\HuangMonitor}[3][1]{%
  \begin{scope}[shift={(#2,#3)}, scale=#1]
    \fill[black] (-0.72,-0.54) rectangle (0.72,0.55);
    \fill[huangLavender!60] (-0.6,-0.4) rectangle (0.6,0.4);
    \fill[white, opacity=0.72] (-0.55,-0.35) -- (-0.12,0.35) -- (0.6,0.35) -- (0.6,-0.35) -- cycle;
    \node[font=\scriptsize\bfseries, text=huangGray!80] at (0.12,0.02) {Model $\Psi$};
    \fill[black] (0,-0.47) circle (0.06);
    \fill[black] (-0.13,-0.8) rectangle (0.13,-0.54);
    \draw[black, line width=1.5pt] (-0.38,-0.84) -- (0.38,-0.84);
  \end{scope}%
}

\providecommand{\HuangMeasurementStack}[3][1]{%
  \begin{scope}[shift={(#2,#3)}, scale=#1]
    \foreach \i in {0,...,6} {
      \node[huang meas] at (0,0.18*\i) {};
      \draw[huangGold!80!black] (-0.08,0.18*\i-0.02) -- (0.08,0.18*\i+0.07);
    }
  \end{scope}%
}

\providecommand{\HuangCircuitGrid}[4][1]{%
  \begin{scope}[shift={(#2,#3)}, scale=#1]
    \foreach \x in {-0.75,-0.45,-0.15,0.15,0.45,0.75}
      \draw[huangGray] (\x,0.9) -- (\x,-1.05);
    \foreach \x/\y/\w/\h in {
      -0.75/0.55/.34/.45,-0.45/0.18/.24/.42,-0.15/0.55/.24/.52,0.15/0.18/.50/.42,0.45/0.55/.34/.45,0.75/0.18/.24/.42,
      -0.75/-0.25/.24/.46,-0.45/-0.63/.52/.42,-0.15/-0.25/.24/.52,0.15/-0.63/.52/.42,0.45/-0.25/.24/.46,0.75/-0.63/.34/.42}
      \node[huang gate, minimum width=\w cm, minimum height=\h cm] at (\x,\y) {};
    \node[font=\scriptsize] at (0,1.12) {$#4$};
  \end{scope}%
}

\providecommand{\HuangNeuralNet}[3][1]{%
  \begin{scope}[shift={(#2,#3)}, scale=#1]
    \node[draw, dashed, rounded corners=8pt, fill=huangPaleGreen, minimum width=1.65cm, minimum height=2.45cm] at (0,0) {};
    \foreach \x in {-0.45,0,0.45} \fill[huangBlue!38] (\x,0.75) circle (0.08);
    \foreach \x in {-0.58,-0.18,0.22,0.58} \fill[huangBlue!38] (\x,0.15) circle (0.08);
    \foreach \x in {-0.58,-0.18,0.22,0.58} \fill[huangBlue!38] (\x,-0.45) circle (0.08);
    \fill[huangBlue!38] (0,-1.02) circle (0.08);
    \foreach \a in {-0.45,0,0.45} \foreach \b in {-0.58,-0.18,0.22,0.58}
      \draw[huangGray!65, thin] (\a,0.75) -- (\b,0.15);
    \foreach \a in {-0.58,-0.18,0.22,0.58} \foreach \b in {-0.58,-0.18,0.22,0.58}
      \draw[huangGray!65, thin] (\a,0.15) -- (\b,-0.45);
    \foreach \a in {-0.58,-0.18,0.22,0.58}
      \draw[huangGray!65, thin] (\a,-0.45) -- (0,-1.02);
  \end{scope}%
}

\providecommand{\HuangDeviceIcon}[3][1]{%
  \begin{scope}[shift={(#2,#3)}, scale=#1]
    \fill[brown!80!black] (-0.72,0.82) rectangle (0.72,1.12);
    \fill[brown!70!black] (-0.82,0.72) rectangle (0.82,0.78);
    \foreach \x in {-0.48,-0.28,0.28,0.48}
      \draw[line width=1.5pt, huangGold] (\x,0.7) -- (\x,-0.55);
    \foreach \y in {0.42,0.12,-0.18}
      \draw[line width=1.2pt, brown!70!black] (-0.65,\y) -- (0.65,\y);
    \foreach \x in {-0.48,-0.38,0.38,0.48}
      \foreach \y in {0.28,-0.08}
        \draw[line width=1pt, huangGold] (\x,\y) circle (0.055);
    \foreach \x in {-0.52,-0.42,0.42,0.52}
      \foreach \y in {-0.38,-0.48,-0.58}
        \draw[line width=1.5pt, huangGold] (\x-0.08,\y) -- (\x+0.08,\y);
    \fill[huangGold!70] (-0.28,-0.6) rectangle (0.28,-0.85);
    \fill[brown!85!black] (-0.15,-0.85) rectangle (0.15,-1.18);
  \end{scope}%
}

\providecommand{\HuangHypercube}[3][1]{%
  \begin{scope}[shift={(#2,#3)}, scale=#1]
    \coordinate (z000) at (-1.2,-0.25);
    \coordinate (z001) at (-0.25,0.85);
    \coordinate (z010) at (-0.25,-0.25);
    \coordinate (z011) at (0.65,0.85);
    \coordinate (z100) at (-0.25,-1.2);
    \coordinate (z101) at (0.65,-0.25);
    \coordinate (z110) at (0.65,-1.2);
    \coordinate (z111) at (1.55,-0.25);
    \foreach \a/\b in {z000/z001,z000/z010,z000/z100,z001/z011,z001/z101,z010/z011,z010/z110,z100/z101,z100/z110,z011/z111,z101/z111,z110/z111}
      \draw[line width=1.5pt, huangGray] (\a) -- (\b);
    \foreach \p/\lab/\anch in {z000/000/below left,z001/001/above,z010/010/above left,z011/011/above,z100/100/below,z101/101/above,z110/110/below,z111/111/right}
      \fill[huangRed] (\p) circle (0.08) node[\anch, font=\scriptsize, text=black] {\lab};
  \end{scope}%
}
}
\begin{tikzpicture}[>=Stealth, font=\footnotesize]
\node[font=\bfseries] at (-0.2,1.65) {Random Clifford};
\HuangLabPatch[0.75]{-1.25}{0.75}
\foreach \y in {0.35,0.55,0.75,0.95,1.15} \draw[huangGray] (-0.33,\y) -- (0.10,\y);
\HuangCircuitGrid[0.72]{1.0}{0.75}{}
\HuangMeasurementStack[0.9]{1.95}{0.2}
\node[align=center, text width=3.3cm] at (0.25,-0.85)
  {Deep random circuits\\and $Z$-basis measurements};

\HuangUseBoundingBox{-2.45}{-1.75}{2.45}{2.0}

\end{tikzpicture}
}}
            \end{center}

            \vspace{0.4em}
            \begin{block}{Random Clifford shadows}
                Number of samples fixed wrt dimension

                Deep circuits and exponential classical postprocessing
            \end{block}
        \end{column}
        \begin{column}<2->{0.48\textwidth}
            \begin{center}
                \makebox[0pt][c]{\hspace*{-2.5cm}\resizebox{0.95\linewidth}{!}{\IfFileExists{../diagrams/huang_templates.tex}{\definecolor{huangPurple}{RGB}{113,49,169}
\definecolor{huangBlue}{RGB}{65,109,190}
\definecolor{huangRed}{RGB}{196,0,0}
\definecolor{huangOrange}{RGB}{239,126,45}
\definecolor{huangGold}{RGB}{196,145,15}
\definecolor{huangGreen}{RGB}{91,164,61}
\definecolor{huangGray}{RGB}{120,120,120}
\definecolor{huangPaleBlue}{RGB}{224,234,238}
\definecolor{huangPaleGreen}{RGB}{225,240,216}
\definecolor{huangPeach}{RGB}{246,185,139}
\definecolor{huangLavender}{RGB}{216,205,236}

\tikzset{
  huang arrow/.style={-Stealth, thick, draw=huangGray},
  huang dashed box/.style={draw, dashed, rounded corners=5pt, fill=huangPaleBlue, inner sep=5pt},
  huang bit/.style={draw, dashed, rounded corners=4pt, minimum width=0.42cm, minimum height=0.36cm, inner sep=1pt, font=\scriptsize\bfseries},
  huang gate/.style={draw, rounded corners=2pt, fill=huangPeach, minimum width=0.35cm, minimum height=0.54cm, inner sep=1pt},
  huang meas/.style={draw, rounded corners=2pt, fill=huangLavender, minimum width=0.24cm, minimum height=0.24cm, inner sep=0pt},
  huang label/.style={font=\scriptsize},
  huang title/.style={font=\bfseries}
}

\providecommand{\HuangAtom}[2]{%
  \begin{scope}[shift={(#1,#2)}, scale=0.18]
    \foreach \a in {0,60,...,300} {
      \draw[huangPurple, thick, rotate=\a] (0,0) ellipse (0.36 and 0.9);
    }
    \fill[huangGreen!85!black] (0,0) circle (0.11);
  \end{scope}%
}

\providecommand{\HuangLabPatch}[3][1]{%
  \begin{scope}[shift={(#2,#3)}, scale=#1]
    \node[huang dashed box, minimum width=2.25cm, minimum height=1.55cm] (labpatch) at (0,0) {};
    \foreach \x/\y in {-0.7/0.43,0/0.43,0.7/0.43,-0.7/-0.05,0/-0.05,0.7/-0.05,-0.35/-0.53,0.35/-0.53}
      {\HuangAtom{\x}{\y}}
  \end{scope}%
}

\providecommand{\HuangLabStack}[3][1]{%
  \begin{scope}[shift={(#2,#3)}, scale=#1]
    \foreach \dx in {-0.42,-0.18,0.06} {
      \node[huang dashed box, minimum width=1.75cm, minimum height=1.25cm, fill=huangPaleBlue!85] at (\dx,0) {};
    }
    \foreach \x/\y in {-0.58/0.36,0/0.36,0.58/0.36,-0.58/-0.03,0/-0.03,0.58/-0.03,-0.3/-0.43,0.3/-0.43}
      {\HuangAtom{\x}{\y}}
  \end{scope}%
}

\providecommand{\HuangMonitor}[3][1]{%
  \begin{scope}[shift={(#2,#3)}, scale=#1]
    \fill[black] (-0.72,-0.54) rectangle (0.72,0.55);
    \fill[huangLavender!60] (-0.6,-0.4) rectangle (0.6,0.4);
    \fill[white, opacity=0.72] (-0.55,-0.35) -- (-0.12,0.35) -- (0.6,0.35) -- (0.6,-0.35) -- cycle;
    \node[font=\scriptsize\bfseries, text=huangGray!80] at (0.12,0.02) {Model $\Psi$};
    \fill[black] (0,-0.47) circle (0.06);
    \fill[black] (-0.13,-0.8) rectangle (0.13,-0.54);
    \draw[black, line width=1.5pt] (-0.38,-0.84) -- (0.38,-0.84);
  \end{scope}%
}

\providecommand{\HuangMeasurementStack}[3][1]{%
  \begin{scope}[shift={(#2,#3)}, scale=#1]
    \foreach \i in {0,...,6} {
      \node[huang meas] at (0,0.18*\i) {};
      \draw[huangGold!80!black] (-0.08,0.18*\i-0.02) -- (0.08,0.18*\i+0.07);
    }
  \end{scope}%
}

\providecommand{\HuangCircuitGrid}[4][1]{%
  \begin{scope}[shift={(#2,#3)}, scale=#1]
    \foreach \x in {-0.75,-0.45,-0.15,0.15,0.45,0.75}
      \draw[huangGray] (\x,0.9) -- (\x,-1.05);
    \foreach \x/\y/\w/\h in {
      -0.75/0.55/.34/.45,-0.45/0.18/.24/.42,-0.15/0.55/.24/.52,0.15/0.18/.50/.42,0.45/0.55/.34/.45,0.75/0.18/.24/.42,
      -0.75/-0.25/.24/.46,-0.45/-0.63/.52/.42,-0.15/-0.25/.24/.52,0.15/-0.63/.52/.42,0.45/-0.25/.24/.46,0.75/-0.63/.34/.42}
      \node[huang gate, minimum width=\w cm, minimum height=\h cm] at (\x,\y) {};
    \node[font=\scriptsize] at (0,1.12) {$#4$};
  \end{scope}%
}

\providecommand{\HuangNeuralNet}[3][1]{%
  \begin{scope}[shift={(#2,#3)}, scale=#1]
    \node[draw, dashed, rounded corners=8pt, fill=huangPaleGreen, minimum width=1.65cm, minimum height=2.45cm] at (0,0) {};
    \foreach \x in {-0.45,0,0.45} \fill[huangBlue!38] (\x,0.75) circle (0.08);
    \foreach \x in {-0.58,-0.18,0.22,0.58} \fill[huangBlue!38] (\x,0.15) circle (0.08);
    \foreach \x in {-0.58,-0.18,0.22,0.58} \fill[huangBlue!38] (\x,-0.45) circle (0.08);
    \fill[huangBlue!38] (0,-1.02) circle (0.08);
    \foreach \a in {-0.45,0,0.45} \foreach \b in {-0.58,-0.18,0.22,0.58}
      \draw[huangGray!65, thin] (\a,0.75) -- (\b,0.15);
    \foreach \a in {-0.58,-0.18,0.22,0.58} \foreach \b in {-0.58,-0.18,0.22,0.58}
      \draw[huangGray!65, thin] (\a,0.15) -- (\b,-0.45);
    \foreach \a in {-0.58,-0.18,0.22,0.58}
      \draw[huangGray!65, thin] (\a,-0.45) -- (0,-1.02);
  \end{scope}%
}

\providecommand{\HuangDeviceIcon}[3][1]{%
  \begin{scope}[shift={(#2,#3)}, scale=#1]
    \fill[brown!80!black] (-0.72,0.82) rectangle (0.72,1.12);
    \fill[brown!70!black] (-0.82,0.72) rectangle (0.82,0.78);
    \foreach \x in {-0.48,-0.28,0.28,0.48}
      \draw[line width=1.5pt, huangGold] (\x,0.7) -- (\x,-0.55);
    \foreach \y in {0.42,0.12,-0.18}
      \draw[line width=1.2pt, brown!70!black] (-0.65,\y) -- (0.65,\y);
    \foreach \x in {-0.48,-0.38,0.38,0.48}
      \foreach \y in {0.28,-0.08}
        \draw[line width=1pt, huangGold] (\x,\y) circle (0.055);
    \foreach \x in {-0.52,-0.42,0.42,0.52}
      \foreach \y in {-0.38,-0.48,-0.58}
        \draw[line width=1.5pt, huangGold] (\x-0.08,\y) -- (\x+0.08,\y);
    \fill[huangGold!70] (-0.28,-0.6) rectangle (0.28,-0.85);
    \fill[brown!85!black] (-0.15,-0.85) rectangle (0.15,-1.18);
  \end{scope}%
}

\providecommand{\HuangHypercube}[3][1]{%
  \begin{scope}[shift={(#2,#3)}, scale=#1]
    \coordinate (z000) at (-1.2,-0.25);
    \coordinate (z001) at (-0.25,0.85);
    \coordinate (z010) at (-0.25,-0.25);
    \coordinate (z011) at (0.65,0.85);
    \coordinate (z100) at (-0.25,-1.2);
    \coordinate (z101) at (0.65,-0.25);
    \coordinate (z110) at (0.65,-1.2);
    \coordinate (z111) at (1.55,-0.25);
    \foreach \a/\b in {z000/z001,z000/z010,z000/z100,z001/z011,z001/z101,z010/z011,z010/z110,z100/z101,z100/z110,z011/z111,z101/z111,z110/z111}
      \draw[line width=1.5pt, huangGray] (\a) -- (\b);
    \foreach \p/\lab/\anch in {z000/000/below left,z001/001/above,z010/010/above left,z011/011/above,z100/100/below,z101/101/above,z110/110/below,z111/111/right}
      \fill[huangRed] (\p) circle (0.08) node[\anch, font=\scriptsize, text=black] {\lab};
  \end{scope}%
}
}{\definecolor{huangPurple}{RGB}{113,49,169}
\definecolor{huangBlue}{RGB}{65,109,190}
\definecolor{huangRed}{RGB}{196,0,0}
\definecolor{huangOrange}{RGB}{239,126,45}
\definecolor{huangGold}{RGB}{196,145,15}
\definecolor{huangGreen}{RGB}{91,164,61}
\definecolor{huangGray}{RGB}{120,120,120}
\definecolor{huangPaleBlue}{RGB}{224,234,238}
\definecolor{huangPaleGreen}{RGB}{225,240,216}
\definecolor{huangPeach}{RGB}{246,185,139}
\definecolor{huangLavender}{RGB}{216,205,236}

\tikzset{
  huang arrow/.style={-Stealth, thick, draw=huangGray},
  huang dashed box/.style={draw, dashed, rounded corners=5pt, fill=huangPaleBlue, inner sep=5pt},
  huang bit/.style={draw, dashed, rounded corners=4pt, minimum width=0.42cm, minimum height=0.36cm, inner sep=1pt, font=\scriptsize\bfseries},
  huang gate/.style={draw, rounded corners=2pt, fill=huangPeach, minimum width=0.35cm, minimum height=0.54cm, inner sep=1pt},
  huang meas/.style={draw, rounded corners=2pt, fill=huangLavender, minimum width=0.24cm, minimum height=0.24cm, inner sep=0pt},
  huang label/.style={font=\scriptsize},
  huang title/.style={font=\bfseries}
}

\providecommand{\HuangAtom}[2]{%
  \begin{scope}[shift={(#1,#2)}, scale=0.18]
    \foreach \a in {0,60,...,300} {
      \draw[huangPurple, thick, rotate=\a] (0,0) ellipse (0.36 and 0.9);
    }
    \fill[huangGreen!85!black] (0,0) circle (0.11);
  \end{scope}%
}

\providecommand{\HuangLabPatch}[3][1]{%
  \begin{scope}[shift={(#2,#3)}, scale=#1]
    \node[huang dashed box, minimum width=2.25cm, minimum height=1.55cm] (labpatch) at (0,0) {};
    \foreach \x/\y in {-0.7/0.43,0/0.43,0.7/0.43,-0.7/-0.05,0/-0.05,0.7/-0.05,-0.35/-0.53,0.35/-0.53}
      {\HuangAtom{\x}{\y}}
  \end{scope}%
}

\providecommand{\HuangLabStack}[3][1]{%
  \begin{scope}[shift={(#2,#3)}, scale=#1]
    \foreach \dx in {-0.42,-0.18,0.06} {
      \node[huang dashed box, minimum width=1.75cm, minimum height=1.25cm, fill=huangPaleBlue!85] at (\dx,0) {};
    }
    \foreach \x/\y in {-0.58/0.36,0/0.36,0.58/0.36,-0.58/-0.03,0/-0.03,0.58/-0.03,-0.3/-0.43,0.3/-0.43}
      {\HuangAtom{\x}{\y}}
  \end{scope}%
}

\providecommand{\HuangMonitor}[3][1]{%
  \begin{scope}[shift={(#2,#3)}, scale=#1]
    \fill[black] (-0.72,-0.54) rectangle (0.72,0.55);
    \fill[huangLavender!60] (-0.6,-0.4) rectangle (0.6,0.4);
    \fill[white, opacity=0.72] (-0.55,-0.35) -- (-0.12,0.35) -- (0.6,0.35) -- (0.6,-0.35) -- cycle;
    \node[font=\scriptsize\bfseries, text=huangGray!80] at (0.12,0.02) {Model $\Psi$};
    \fill[black] (0,-0.47) circle (0.06);
    \fill[black] (-0.13,-0.8) rectangle (0.13,-0.54);
    \draw[black, line width=1.5pt] (-0.38,-0.84) -- (0.38,-0.84);
  \end{scope}%
}

\providecommand{\HuangMeasurementStack}[3][1]{%
  \begin{scope}[shift={(#2,#3)}, scale=#1]
    \foreach \i in {0,...,6} {
      \node[huang meas] at (0,0.18*\i) {};
      \draw[huangGold!80!black] (-0.08,0.18*\i-0.02) -- (0.08,0.18*\i+0.07);
    }
  \end{scope}%
}

\providecommand{\HuangCircuitGrid}[4][1]{%
  \begin{scope}[shift={(#2,#3)}, scale=#1]
    \foreach \x in {-0.75,-0.45,-0.15,0.15,0.45,0.75}
      \draw[huangGray] (\x,0.9) -- (\x,-1.05);
    \foreach \x/\y/\w/\h in {
      -0.75/0.55/.34/.45,-0.45/0.18/.24/.42,-0.15/0.55/.24/.52,0.15/0.18/.50/.42,0.45/0.55/.34/.45,0.75/0.18/.24/.42,
      -0.75/-0.25/.24/.46,-0.45/-0.63/.52/.42,-0.15/-0.25/.24/.52,0.15/-0.63/.52/.42,0.45/-0.25/.24/.46,0.75/-0.63/.34/.42}
      \node[huang gate, minimum width=\w cm, minimum height=\h cm] at (\x,\y) {};
    \node[font=\scriptsize] at (0,1.12) {$#4$};
  \end{scope}%
}

\providecommand{\HuangNeuralNet}[3][1]{%
  \begin{scope}[shift={(#2,#3)}, scale=#1]
    \node[draw, dashed, rounded corners=8pt, fill=huangPaleGreen, minimum width=1.65cm, minimum height=2.45cm] at (0,0) {};
    \foreach \x in {-0.45,0,0.45} \fill[huangBlue!38] (\x,0.75) circle (0.08);
    \foreach \x in {-0.58,-0.18,0.22,0.58} \fill[huangBlue!38] (\x,0.15) circle (0.08);
    \foreach \x in {-0.58,-0.18,0.22,0.58} \fill[huangBlue!38] (\x,-0.45) circle (0.08);
    \fill[huangBlue!38] (0,-1.02) circle (0.08);
    \foreach \a in {-0.45,0,0.45} \foreach \b in {-0.58,-0.18,0.22,0.58}
      \draw[huangGray!65, thin] (\a,0.75) -- (\b,0.15);
    \foreach \a in {-0.58,-0.18,0.22,0.58} \foreach \b in {-0.58,-0.18,0.22,0.58}
      \draw[huangGray!65, thin] (\a,0.15) -- (\b,-0.45);
    \foreach \a in {-0.58,-0.18,0.22,0.58}
      \draw[huangGray!65, thin] (\a,-0.45) -- (0,-1.02);
  \end{scope}%
}

\providecommand{\HuangDeviceIcon}[3][1]{%
  \begin{scope}[shift={(#2,#3)}, scale=#1]
    \fill[brown!80!black] (-0.72,0.82) rectangle (0.72,1.12);
    \fill[brown!70!black] (-0.82,0.72) rectangle (0.82,0.78);
    \foreach \x in {-0.48,-0.28,0.28,0.48}
      \draw[line width=1.5pt, huangGold] (\x,0.7) -- (\x,-0.55);
    \foreach \y in {0.42,0.12,-0.18}
      \draw[line width=1.2pt, brown!70!black] (-0.65,\y) -- (0.65,\y);
    \foreach \x in {-0.48,-0.38,0.38,0.48}
      \foreach \y in {0.28,-0.08}
        \draw[line width=1pt, huangGold] (\x,\y) circle (0.055);
    \foreach \x in {-0.52,-0.42,0.42,0.52}
      \foreach \y in {-0.38,-0.48,-0.58}
        \draw[line width=1.5pt, huangGold] (\x-0.08,\y) -- (\x+0.08,\y);
    \fill[huangGold!70] (-0.28,-0.6) rectangle (0.28,-0.85);
    \fill[brown!85!black] (-0.15,-0.85) rectangle (0.15,-1.18);
  \end{scope}%
}

\providecommand{\HuangHypercube}[3][1]{%
  \begin{scope}[shift={(#2,#3)}, scale=#1]
    \coordinate (z000) at (-1.2,-0.25);
    \coordinate (z001) at (-0.25,0.85);
    \coordinate (z010) at (-0.25,-0.25);
    \coordinate (z011) at (0.65,0.85);
    \coordinate (z100) at (-0.25,-1.2);
    \coordinate (z101) at (0.65,-0.25);
    \coordinate (z110) at (0.65,-1.2);
    \coordinate (z111) at (1.55,-0.25);
    \foreach \a/\b in {z000/z001,z000/z010,z000/z100,z001/z011,z001/z101,z010/z011,z010/z110,z100/z101,z100/z110,z011/z111,z101/z111,z110/z111}
      \draw[line width=1.5pt, huangGray] (\a) -- (\b);
    \foreach \p/\lab/\anch in {z000/000/below left,z001/001/above,z010/010/above left,z011/011/above,z100/100/below,z101/101/above,z110/110/below,z111/111/right}
      \fill[huangRed] (\p) circle (0.08) node[\anch, font=\scriptsize, text=black] {\lab};
  \end{scope}%
}
}
\begin{tikzpicture}[>=Stealth, font=\footnotesize]
\node[font=\bfseries] at (0,1.65) {Random Pauli};
\node[draw, rounded corners=10pt, fill=huangPaleBlue, minimum width=1.85cm, minimum height=1.55cm] at (0,0.75) {};
\foreach \x/\y in {-0.45/1.15,0.05/1.15,-0.45/0.72,0.05/0.72,0.5/0.72,-0.25/0.3,0.25/0.3}
  {\node[huang bit, fill=orange!16] at (\x,\y) {}; \HuangAtom{\x}{\y}}
\node[huang bit, fill=orange!18] at (0.85,-0.15) {};
\node[text=huangOrange] at (0,-0.65) {Random X/Y/Z-basis};
\node[align=center, text width=3.1cm] at (0,-1.15)
  {Single-qubit Pauli bases,\\but exponentially many samples};

\HuangUseBoundingBox{-2.45}{-1.75}{2.45}{2.0}

\end{tikzpicture}
}}
            \end{center}

            \vspace{0.4em}
            \begin{block}{Single-qubit Pauli shadows}
                Number of samples grows exponentially

                Simple to build
            \end{block}
        \end{column}
    \end{columns}
\end{frame}

\begin{frame}{Restricted Measurements}
    \begin{itemize}
        \item Direct fidelity estimation and Pauli-shadow variants work well for structured observables and target families.
        \item Efficient certification is known for stabilizer states, shallow-circuit outputs, hypergraph states, IQP outputs, and Gaussian states.
        \item For generic random targets, XEB can work under suitable noise assumptions, but it is not a general certification guarantee.
    \end{itemize}

    \vspace{0.4em}
    \begin{center}
        \resizebox{0.58\textwidth}{!}{\IfFileExists{../diagrams/huang_templates.tex}{\definecolor{huangPurple}{RGB}{113,49,169}
\definecolor{huangBlue}{RGB}{65,109,190}
\definecolor{huangRed}{RGB}{196,0,0}
\definecolor{huangOrange}{RGB}{239,126,45}
\definecolor{huangGold}{RGB}{196,145,15}
\definecolor{huangGreen}{RGB}{91,164,61}
\definecolor{huangGray}{RGB}{120,120,120}
\definecolor{huangPaleBlue}{RGB}{224,234,238}
\definecolor{huangPaleGreen}{RGB}{225,240,216}
\definecolor{huangPeach}{RGB}{246,185,139}
\definecolor{huangLavender}{RGB}{216,205,236}

\tikzset{
  huang arrow/.style={-Stealth, thick, draw=huangGray},
  huang dashed box/.style={draw, dashed, rounded corners=5pt, fill=huangPaleBlue, inner sep=5pt},
  huang bit/.style={draw, dashed, rounded corners=4pt, minimum width=0.42cm, minimum height=0.36cm, inner sep=1pt, font=\scriptsize\bfseries},
  huang gate/.style={draw, rounded corners=2pt, fill=huangPeach, minimum width=0.35cm, minimum height=0.54cm, inner sep=1pt},
  huang meas/.style={draw, rounded corners=2pt, fill=huangLavender, minimum width=0.24cm, minimum height=0.24cm, inner sep=0pt},
  huang label/.style={font=\scriptsize},
  huang title/.style={font=\bfseries}
}

\providecommand{\HuangAtom}[2]{%
  \begin{scope}[shift={(#1,#2)}, scale=0.18]
    \foreach \a in {0,60,...,300} {
      \draw[huangPurple, thick, rotate=\a] (0,0) ellipse (0.36 and 0.9);
    }
    \fill[huangGreen!85!black] (0,0) circle (0.11);
  \end{scope}%
}

\providecommand{\HuangLabPatch}[3][1]{%
  \begin{scope}[shift={(#2,#3)}, scale=#1]
    \node[huang dashed box, minimum width=2.25cm, minimum height=1.55cm] (labpatch) at (0,0) {};
    \foreach \x/\y in {-0.7/0.43,0/0.43,0.7/0.43,-0.7/-0.05,0/-0.05,0.7/-0.05,-0.35/-0.53,0.35/-0.53}
      {\HuangAtom{\x}{\y}}
  \end{scope}%
}

\providecommand{\HuangLabStack}[3][1]{%
  \begin{scope}[shift={(#2,#3)}, scale=#1]
    \foreach \dx in {-0.42,-0.18,0.06} {
      \node[huang dashed box, minimum width=1.75cm, minimum height=1.25cm, fill=huangPaleBlue!85] at (\dx,0) {};
    }
    \foreach \x/\y in {-0.58/0.36,0/0.36,0.58/0.36,-0.58/-0.03,0/-0.03,0.58/-0.03,-0.3/-0.43,0.3/-0.43}
      {\HuangAtom{\x}{\y}}
  \end{scope}%
}

\providecommand{\HuangMonitor}[3][1]{%
  \begin{scope}[shift={(#2,#3)}, scale=#1]
    \fill[black] (-0.72,-0.54) rectangle (0.72,0.55);
    \fill[huangLavender!60] (-0.6,-0.4) rectangle (0.6,0.4);
    \fill[white, opacity=0.72] (-0.55,-0.35) -- (-0.12,0.35) -- (0.6,0.35) -- (0.6,-0.35) -- cycle;
    \node[font=\scriptsize\bfseries, text=huangGray!80] at (0.12,0.02) {Model $\Psi$};
    \fill[black] (0,-0.47) circle (0.06);
    \fill[black] (-0.13,-0.8) rectangle (0.13,-0.54);
    \draw[black, line width=1.5pt] (-0.38,-0.84) -- (0.38,-0.84);
  \end{scope}%
}

\providecommand{\HuangMeasurementStack}[3][1]{%
  \begin{scope}[shift={(#2,#3)}, scale=#1]
    \foreach \i in {0,...,6} {
      \node[huang meas] at (0,0.18*\i) {};
      \draw[huangGold!80!black] (-0.08,0.18*\i-0.02) -- (0.08,0.18*\i+0.07);
    }
  \end{scope}%
}

\providecommand{\HuangCircuitGrid}[4][1]{%
  \begin{scope}[shift={(#2,#3)}, scale=#1]
    \foreach \x in {-0.75,-0.45,-0.15,0.15,0.45,0.75}
      \draw[huangGray] (\x,0.9) -- (\x,-1.05);
    \foreach \x/\y/\w/\h in {
      -0.75/0.55/.34/.45,-0.45/0.18/.24/.42,-0.15/0.55/.24/.52,0.15/0.18/.50/.42,0.45/0.55/.34/.45,0.75/0.18/.24/.42,
      -0.75/-0.25/.24/.46,-0.45/-0.63/.52/.42,-0.15/-0.25/.24/.52,0.15/-0.63/.52/.42,0.45/-0.25/.24/.46,0.75/-0.63/.34/.42}
      \node[huang gate, minimum width=\w cm, minimum height=\h cm] at (\x,\y) {};
    \node[font=\scriptsize] at (0,1.12) {$#4$};
  \end{scope}%
}

\providecommand{\HuangNeuralNet}[3][1]{%
  \begin{scope}[shift={(#2,#3)}, scale=#1]
    \node[draw, dashed, rounded corners=8pt, fill=huangPaleGreen, minimum width=1.65cm, minimum height=2.45cm] at (0,0) {};
    \foreach \x in {-0.45,0,0.45} \fill[huangBlue!38] (\x,0.75) circle (0.08);
    \foreach \x in {-0.58,-0.18,0.22,0.58} \fill[huangBlue!38] (\x,0.15) circle (0.08);
    \foreach \x in {-0.58,-0.18,0.22,0.58} \fill[huangBlue!38] (\x,-0.45) circle (0.08);
    \fill[huangBlue!38] (0,-1.02) circle (0.08);
    \foreach \a in {-0.45,0,0.45} \foreach \b in {-0.58,-0.18,0.22,0.58}
      \draw[huangGray!65, thin] (\a,0.75) -- (\b,0.15);
    \foreach \a in {-0.58,-0.18,0.22,0.58} \foreach \b in {-0.58,-0.18,0.22,0.58}
      \draw[huangGray!65, thin] (\a,0.15) -- (\b,-0.45);
    \foreach \a in {-0.58,-0.18,0.22,0.58}
      \draw[huangGray!65, thin] (\a,-0.45) -- (0,-1.02);
  \end{scope}%
}

\providecommand{\HuangDeviceIcon}[3][1]{%
  \begin{scope}[shift={(#2,#3)}, scale=#1]
    \fill[brown!80!black] (-0.72,0.82) rectangle (0.72,1.12);
    \fill[brown!70!black] (-0.82,0.72) rectangle (0.82,0.78);
    \foreach \x in {-0.48,-0.28,0.28,0.48}
      \draw[line width=1.5pt, huangGold] (\x,0.7) -- (\x,-0.55);
    \foreach \y in {0.42,0.12,-0.18}
      \draw[line width=1.2pt, brown!70!black] (-0.65,\y) -- (0.65,\y);
    \foreach \x in {-0.48,-0.38,0.38,0.48}
      \foreach \y in {0.28,-0.08}
        \draw[line width=1pt, huangGold] (\x,\y) circle (0.055);
    \foreach \x in {-0.52,-0.42,0.42,0.52}
      \foreach \y in {-0.38,-0.48,-0.58}
        \draw[line width=1.5pt, huangGold] (\x-0.08,\y) -- (\x+0.08,\y);
    \fill[huangGold!70] (-0.28,-0.6) rectangle (0.28,-0.85);
    \fill[brown!85!black] (-0.15,-0.85) rectangle (0.15,-1.18);
  \end{scope}%
}

\providecommand{\HuangHypercube}[3][1]{%
  \begin{scope}[shift={(#2,#3)}, scale=#1]
    \coordinate (z000) at (-1.2,-0.25);
    \coordinate (z001) at (-0.25,0.85);
    \coordinate (z010) at (-0.25,-0.25);
    \coordinate (z011) at (0.65,0.85);
    \coordinate (z100) at (-0.25,-1.2);
    \coordinate (z101) at (0.65,-0.25);
    \coordinate (z110) at (0.65,-1.2);
    \coordinate (z111) at (1.55,-0.25);
    \foreach \a/\b in {z000/z001,z000/z010,z000/z100,z001/z011,z001/z101,z010/z011,z010/z110,z100/z101,z100/z110,z011/z111,z101/z111,z110/z111}
      \draw[line width=1.5pt, huangGray] (\a) -- (\b);
    \foreach \p/\lab/\anch in {z000/000/below left,z001/001/above,z010/010/above left,z011/011/above,z100/100/below,z101/101/above,z110/110/below,z111/111/right}
      \fill[huangRed] (\p) circle (0.08) node[\anch, font=\scriptsize, text=black] {\lab};
  \end{scope}%
}
}{\definecolor{huangPurple}{RGB}{113,49,169}
\definecolor{huangBlue}{RGB}{65,109,190}
\definecolor{huangRed}{RGB}{196,0,0}
\definecolor{huangOrange}{RGB}{239,126,45}
\definecolor{huangGold}{RGB}{196,145,15}
\definecolor{huangGreen}{RGB}{91,164,61}
\definecolor{huangGray}{RGB}{120,120,120}
\definecolor{huangPaleBlue}{RGB}{224,234,238}
\definecolor{huangPaleGreen}{RGB}{225,240,216}
\definecolor{huangPeach}{RGB}{246,185,139}
\definecolor{huangLavender}{RGB}{216,205,236}

\tikzset{
  huang arrow/.style={-Stealth, thick, draw=huangGray},
  huang dashed box/.style={draw, dashed, rounded corners=5pt, fill=huangPaleBlue, inner sep=5pt},
  huang bit/.style={draw, dashed, rounded corners=4pt, minimum width=0.42cm, minimum height=0.36cm, inner sep=1pt, font=\scriptsize\bfseries},
  huang gate/.style={draw, rounded corners=2pt, fill=huangPeach, minimum width=0.35cm, minimum height=0.54cm, inner sep=1pt},
  huang meas/.style={draw, rounded corners=2pt, fill=huangLavender, minimum width=0.24cm, minimum height=0.24cm, inner sep=0pt},
  huang label/.style={font=\scriptsize},
  huang title/.style={font=\bfseries}
}

\providecommand{\HuangAtom}[2]{%
  \begin{scope}[shift={(#1,#2)}, scale=0.18]
    \foreach \a in {0,60,...,300} {
      \draw[huangPurple, thick, rotate=\a] (0,0) ellipse (0.36 and 0.9);
    }
    \fill[huangGreen!85!black] (0,0) circle (0.11);
  \end{scope}%
}

\providecommand{\HuangLabPatch}[3][1]{%
  \begin{scope}[shift={(#2,#3)}, scale=#1]
    \node[huang dashed box, minimum width=2.25cm, minimum height=1.55cm] (labpatch) at (0,0) {};
    \foreach \x/\y in {-0.7/0.43,0/0.43,0.7/0.43,-0.7/-0.05,0/-0.05,0.7/-0.05,-0.35/-0.53,0.35/-0.53}
      {\HuangAtom{\x}{\y}}
  \end{scope}%
}

\providecommand{\HuangLabStack}[3][1]{%
  \begin{scope}[shift={(#2,#3)}, scale=#1]
    \foreach \dx in {-0.42,-0.18,0.06} {
      \node[huang dashed box, minimum width=1.75cm, minimum height=1.25cm, fill=huangPaleBlue!85] at (\dx,0) {};
    }
    \foreach \x/\y in {-0.58/0.36,0/0.36,0.58/0.36,-0.58/-0.03,0/-0.03,0.58/-0.03,-0.3/-0.43,0.3/-0.43}
      {\HuangAtom{\x}{\y}}
  \end{scope}%
}

\providecommand{\HuangMonitor}[3][1]{%
  \begin{scope}[shift={(#2,#3)}, scale=#1]
    \fill[black] (-0.72,-0.54) rectangle (0.72,0.55);
    \fill[huangLavender!60] (-0.6,-0.4) rectangle (0.6,0.4);
    \fill[white, opacity=0.72] (-0.55,-0.35) -- (-0.12,0.35) -- (0.6,0.35) -- (0.6,-0.35) -- cycle;
    \node[font=\scriptsize\bfseries, text=huangGray!80] at (0.12,0.02) {Model $\Psi$};
    \fill[black] (0,-0.47) circle (0.06);
    \fill[black] (-0.13,-0.8) rectangle (0.13,-0.54);
    \draw[black, line width=1.5pt] (-0.38,-0.84) -- (0.38,-0.84);
  \end{scope}%
}

\providecommand{\HuangMeasurementStack}[3][1]{%
  \begin{scope}[shift={(#2,#3)}, scale=#1]
    \foreach \i in {0,...,6} {
      \node[huang meas] at (0,0.18*\i) {};
      \draw[huangGold!80!black] (-0.08,0.18*\i-0.02) -- (0.08,0.18*\i+0.07);
    }
  \end{scope}%
}

\providecommand{\HuangCircuitGrid}[4][1]{%
  \begin{scope}[shift={(#2,#3)}, scale=#1]
    \foreach \x in {-0.75,-0.45,-0.15,0.15,0.45,0.75}
      \draw[huangGray] (\x,0.9) -- (\x,-1.05);
    \foreach \x/\y/\w/\h in {
      -0.75/0.55/.34/.45,-0.45/0.18/.24/.42,-0.15/0.55/.24/.52,0.15/0.18/.50/.42,0.45/0.55/.34/.45,0.75/0.18/.24/.42,
      -0.75/-0.25/.24/.46,-0.45/-0.63/.52/.42,-0.15/-0.25/.24/.52,0.15/-0.63/.52/.42,0.45/-0.25/.24/.46,0.75/-0.63/.34/.42}
      \node[huang gate, minimum width=\w cm, minimum height=\h cm] at (\x,\y) {};
    \node[font=\scriptsize] at (0,1.12) {$#4$};
  \end{scope}%
}

\providecommand{\HuangNeuralNet}[3][1]{%
  \begin{scope}[shift={(#2,#3)}, scale=#1]
    \node[draw, dashed, rounded corners=8pt, fill=huangPaleGreen, minimum width=1.65cm, minimum height=2.45cm] at (0,0) {};
    \foreach \x in {-0.45,0,0.45} \fill[huangBlue!38] (\x,0.75) circle (0.08);
    \foreach \x in {-0.58,-0.18,0.22,0.58} \fill[huangBlue!38] (\x,0.15) circle (0.08);
    \foreach \x in {-0.58,-0.18,0.22,0.58} \fill[huangBlue!38] (\x,-0.45) circle (0.08);
    \fill[huangBlue!38] (0,-1.02) circle (0.08);
    \foreach \a in {-0.45,0,0.45} \foreach \b in {-0.58,-0.18,0.22,0.58}
      \draw[huangGray!65, thin] (\a,0.75) -- (\b,0.15);
    \foreach \a in {-0.58,-0.18,0.22,0.58} \foreach \b in {-0.58,-0.18,0.22,0.58}
      \draw[huangGray!65, thin] (\a,0.15) -- (\b,-0.45);
    \foreach \a in {-0.58,-0.18,0.22,0.58}
      \draw[huangGray!65, thin] (\a,-0.45) -- (0,-1.02);
  \end{scope}%
}

\providecommand{\HuangDeviceIcon}[3][1]{%
  \begin{scope}[shift={(#2,#3)}, scale=#1]
    \fill[brown!80!black] (-0.72,0.82) rectangle (0.72,1.12);
    \fill[brown!70!black] (-0.82,0.72) rectangle (0.82,0.78);
    \foreach \x in {-0.48,-0.28,0.28,0.48}
      \draw[line width=1.5pt, huangGold] (\x,0.7) -- (\x,-0.55);
    \foreach \y in {0.42,0.12,-0.18}
      \draw[line width=1.2pt, brown!70!black] (-0.65,\y) -- (0.65,\y);
    \foreach \x in {-0.48,-0.38,0.38,0.48}
      \foreach \y in {0.28,-0.08}
        \draw[line width=1pt, huangGold] (\x,\y) circle (0.055);
    \foreach \x in {-0.52,-0.42,0.42,0.52}
      \foreach \y in {-0.38,-0.48,-0.58}
        \draw[line width=1.5pt, huangGold] (\x-0.08,\y) -- (\x+0.08,\y);
    \fill[huangGold!70] (-0.28,-0.6) rectangle (0.28,-0.85);
    \fill[brown!85!black] (-0.15,-0.85) rectangle (0.15,-1.18);
  \end{scope}%
}

\providecommand{\HuangHypercube}[3][1]{%
  \begin{scope}[shift={(#2,#3)}, scale=#1]
    \coordinate (z000) at (-1.2,-0.25);
    \coordinate (z001) at (-0.25,0.85);
    \coordinate (z010) at (-0.25,-0.25);
    \coordinate (z011) at (0.65,0.85);
    \coordinate (z100) at (-0.25,-1.2);
    \coordinate (z101) at (0.65,-0.25);
    \coordinate (z110) at (0.65,-1.2);
    \coordinate (z111) at (1.55,-0.25);
    \foreach \a/\b in {z000/z001,z000/z010,z000/z100,z001/z011,z001/z101,z010/z011,z010/z110,z100/z101,z100/z110,z011/z111,z101/z111,z110/z111}
      \draw[line width=1.5pt, huangGray] (\a) -- (\b);
    \foreach \p/\lab/\anch in {z000/000/below left,z001/001/above,z010/010/above left,z011/011/above,z100/100/below,z101/101/above,z110/110/below,z111/111/right}
      \fill[huangRed] (\p) circle (0.08) node[\anch, font=\scriptsize, text=black] {\lab};
  \end{scope}%
}
}
\begin{tikzpicture}[>=Stealth, font=\footnotesize]
\node[font=\bfseries] at (0,1.75) {Cross-Entropy Benchmarking};
\node[font=\large] at (0,1.2)
  {$\mathrm{XEB}\propto \mathbb{E}_{z\sim\rho}|\langle z|\psi\rangle|^2-2^{-n}$};
\node[text=huangPurple, align=center] at (0,0.6)
  {Heuristic; can certify incorrectly when $\langle\psi|\rho|\psi\rangle\ll1$};

\node[font=\bfseries] at (0,-0.05) {Lab state $\rho$};
\node[draw, rounded corners=10pt, fill=huangPaleBlue, minimum width=1.85cm, minimum height=1.45cm] at (0,-0.85) {};
\foreach \x/\y in {-0.45/-0.45,0.15/-0.45,-0.65/-0.86,-0.05/-0.86,0.55/-0.86,-0.25/-1.25,0.35/-1.25}
  {\node[huang bit, fill=huangLavender!60] at (\x,\y) {}; \HuangAtom{\x}{\y}}
\node[huang bit, fill=huangLavender!65] at (1.25,-1.45) {};
\node[text=huangPurple] at (0,-1.75) {$Z$-basis};

\HuangUseBoundingBox{-3.25}{-2.0}{3.25}{2.05}

\end{tikzpicture}
}
    \end{center}
\end{frame}

\begin{frame}{Main Theorem}
    \begin{theorem}
      There exists a certification algorithm that takes a classical description of an $n$-qubit {\bf pure} state $\ket{\psi_{\text{target}}}$ and a single copy of an $n$-qubit {\bf mixed} state $\rho_{\text{lab}}$, performs only {\bf single-qubit measurements}, and outputs:
        \begin{itemize}
            \item \textsc{Accept} with probability at least
            $\mathrm{Fid} := \bra{\psi_{\text{target}}}\rho_{\text{lab}}\ket{\psi_{\text{target}}}$
            \item \textsc{Reject} with probability at least
            $\mathrm{Infid}/n := (1-\mathrm{Fid})/n$
        \end{itemize}
    \end{theorem}

    \pause
    To get an arbitrarily high success rate $1 - \epsilon$, just repeat $\frac n {\mathrm{Infid}} \log(1/\epsilon)$ times.
    %\begin{center}
    %    One qubit at a time: compare locally, guarantee globally.
    %\end{center}
\end{frame}

\begin{frame}{Motivating Applications}
    \begin{center}
        \resizebox{0.86\textwidth}{!}{\IfFileExists{../diagrams/huang_templates.tex}{\definecolor{huangPurple}{RGB}{113,49,169}
\definecolor{huangBlue}{RGB}{65,109,190}
\definecolor{huangRed}{RGB}{196,0,0}
\definecolor{huangOrange}{RGB}{239,126,45}
\definecolor{huangGold}{RGB}{196,145,15}
\definecolor{huangGreen}{RGB}{91,164,61}
\definecolor{huangGray}{RGB}{120,120,120}
\definecolor{huangPaleBlue}{RGB}{224,234,238}
\definecolor{huangPaleGreen}{RGB}{225,240,216}
\definecolor{huangPeach}{RGB}{246,185,139}
\definecolor{huangLavender}{RGB}{216,205,236}

\tikzset{
  huang arrow/.style={-Stealth, thick, draw=huangGray},
  huang dashed box/.style={draw, dashed, rounded corners=5pt, fill=huangPaleBlue, inner sep=5pt},
  huang bit/.style={draw, dashed, rounded corners=4pt, minimum width=0.42cm, minimum height=0.36cm, inner sep=1pt, font=\scriptsize\bfseries},
  huang gate/.style={draw, rounded corners=2pt, fill=huangPeach, minimum width=0.35cm, minimum height=0.54cm, inner sep=1pt},
  huang meas/.style={draw, rounded corners=2pt, fill=huangLavender, minimum width=0.24cm, minimum height=0.24cm, inner sep=0pt},
  huang label/.style={font=\scriptsize},
  huang title/.style={font=\bfseries}
}

\providecommand{\HuangAtom}[2]{%
  \begin{scope}[shift={(#1,#2)}, scale=0.18]
    \foreach \a in {0,60,...,300} {
      \draw[huangPurple, thick, rotate=\a] (0,0) ellipse (0.36 and 0.9);
    }
    \fill[huangGreen!85!black] (0,0) circle (0.11);
  \end{scope}%
}

\providecommand{\HuangLabPatch}[3][1]{%
  \begin{scope}[shift={(#2,#3)}, scale=#1]
    \node[huang dashed box, minimum width=2.25cm, minimum height=1.55cm] (labpatch) at (0,0) {};
    \foreach \x/\y in {-0.7/0.43,0/0.43,0.7/0.43,-0.7/-0.05,0/-0.05,0.7/-0.05,-0.35/-0.53,0.35/-0.53}
      {\HuangAtom{\x}{\y}}
  \end{scope}%
}

\providecommand{\HuangLabStack}[3][1]{%
  \begin{scope}[shift={(#2,#3)}, scale=#1]
    \foreach \dx in {-0.42,-0.18,0.06} {
      \node[huang dashed box, minimum width=1.75cm, minimum height=1.25cm, fill=huangPaleBlue!85] at (\dx,0) {};
    }
    \foreach \x/\y in {-0.58/0.36,0/0.36,0.58/0.36,-0.58/-0.03,0/-0.03,0.58/-0.03,-0.3/-0.43,0.3/-0.43}
      {\HuangAtom{\x}{\y}}
  \end{scope}%
}

\providecommand{\HuangMonitor}[3][1]{%
  \begin{scope}[shift={(#2,#3)}, scale=#1]
    \fill[black] (-0.72,-0.54) rectangle (0.72,0.55);
    \fill[huangLavender!60] (-0.6,-0.4) rectangle (0.6,0.4);
    \fill[white, opacity=0.72] (-0.55,-0.35) -- (-0.12,0.35) -- (0.6,0.35) -- (0.6,-0.35) -- cycle;
    \node[font=\scriptsize\bfseries, text=huangGray!80] at (0.12,0.02) {Model $\Psi$};
    \fill[black] (0,-0.47) circle (0.06);
    \fill[black] (-0.13,-0.8) rectangle (0.13,-0.54);
    \draw[black, line width=1.5pt] (-0.38,-0.84) -- (0.38,-0.84);
  \end{scope}%
}

\providecommand{\HuangMeasurementStack}[3][1]{%
  \begin{scope}[shift={(#2,#3)}, scale=#1]
    \foreach \i in {0,...,6} {
      \node[huang meas] at (0,0.18*\i) {};
      \draw[huangGold!80!black] (-0.08,0.18*\i-0.02) -- (0.08,0.18*\i+0.07);
    }
  \end{scope}%
}

\providecommand{\HuangCircuitGrid}[4][1]{%
  \begin{scope}[shift={(#2,#3)}, scale=#1]
    \foreach \x in {-0.75,-0.45,-0.15,0.15,0.45,0.75}
      \draw[huangGray] (\x,0.9) -- (\x,-1.05);
    \foreach \x/\y/\w/\h in {
      -0.75/0.55/.34/.45,-0.45/0.18/.24/.42,-0.15/0.55/.24/.52,0.15/0.18/.50/.42,0.45/0.55/.34/.45,0.75/0.18/.24/.42,
      -0.75/-0.25/.24/.46,-0.45/-0.63/.52/.42,-0.15/-0.25/.24/.52,0.15/-0.63/.52/.42,0.45/-0.25/.24/.46,0.75/-0.63/.34/.42}
      \node[huang gate, minimum width=\w cm, minimum height=\h cm] at (\x,\y) {};
    \node[font=\scriptsize] at (0,1.12) {$#4$};
  \end{scope}%
}

\providecommand{\HuangNeuralNet}[3][1]{%
  \begin{scope}[shift={(#2,#3)}, scale=#1]
    \node[draw, dashed, rounded corners=8pt, fill=huangPaleGreen, minimum width=1.65cm, minimum height=2.45cm] at (0,0) {};
    \foreach \x in {-0.45,0,0.45} \fill[huangBlue!38] (\x,0.75) circle (0.08);
    \foreach \x in {-0.58,-0.18,0.22,0.58} \fill[huangBlue!38] (\x,0.15) circle (0.08);
    \foreach \x in {-0.58,-0.18,0.22,0.58} \fill[huangBlue!38] (\x,-0.45) circle (0.08);
    \fill[huangBlue!38] (0,-1.02) circle (0.08);
    \foreach \a in {-0.45,0,0.45} \foreach \b in {-0.58,-0.18,0.22,0.58}
      \draw[huangGray!65, thin] (\a,0.75) -- (\b,0.15);
    \foreach \a in {-0.58,-0.18,0.22,0.58} \foreach \b in {-0.58,-0.18,0.22,0.58}
      \draw[huangGray!65, thin] (\a,0.15) -- (\b,-0.45);
    \foreach \a in {-0.58,-0.18,0.22,0.58}
      \draw[huangGray!65, thin] (\a,-0.45) -- (0,-1.02);
  \end{scope}%
}

\providecommand{\HuangDeviceIcon}[3][1]{%
  \begin{scope}[shift={(#2,#3)}, scale=#1]
    \fill[brown!80!black] (-0.72,0.82) rectangle (0.72,1.12);
    \fill[brown!70!black] (-0.82,0.72) rectangle (0.82,0.78);
    \foreach \x in {-0.48,-0.28,0.28,0.48}
      \draw[line width=1.5pt, huangGold] (\x,0.7) -- (\x,-0.55);
    \foreach \y in {0.42,0.12,-0.18}
      \draw[line width=1.2pt, brown!70!black] (-0.65,\y) -- (0.65,\y);
    \foreach \x in {-0.48,-0.38,0.38,0.48}
      \foreach \y in {0.28,-0.08}
        \draw[line width=1pt, huangGold] (\x,\y) circle (0.055);
    \foreach \x in {-0.52,-0.42,0.42,0.52}
      \foreach \y in {-0.38,-0.48,-0.58}
        \draw[line width=1.5pt, huangGold] (\x-0.08,\y) -- (\x+0.08,\y);
    \fill[huangGold!70] (-0.28,-0.6) rectangle (0.28,-0.85);
    \fill[brown!85!black] (-0.15,-0.85) rectangle (0.15,-1.18);
  \end{scope}%
}

\providecommand{\HuangHypercube}[3][1]{%
  \begin{scope}[shift={(#2,#3)}, scale=#1]
    \coordinate (z000) at (-1.2,-0.25);
    \coordinate (z001) at (-0.25,0.85);
    \coordinate (z010) at (-0.25,-0.25);
    \coordinate (z011) at (0.65,0.85);
    \coordinate (z100) at (-0.25,-1.2);
    \coordinate (z101) at (0.65,-0.25);
    \coordinate (z110) at (0.65,-1.2);
    \coordinate (z111) at (1.55,-0.25);
    \foreach \a/\b in {z000/z001,z000/z010,z000/z100,z001/z011,z001/z101,z010/z011,z010/z110,z100/z101,z100/z110,z011/z111,z101/z111,z110/z111}
      \draw[line width=1.5pt, huangGray] (\a) -- (\b);
    \foreach \p/\lab/\anch in {z000/000/below left,z001/001/above,z010/010/above left,z011/011/above,z100/100/below,z101/101/above,z110/110/below,z111/111/right}
      \fill[huangRed] (\p) circle (0.08) node[\anch, font=\scriptsize, text=black] {\lab};
  \end{scope}%
}
}{\definecolor{huangPurple}{RGB}{113,49,169}
\definecolor{huangBlue}{RGB}{65,109,190}
\definecolor{huangRed}{RGB}{196,0,0}
\definecolor{huangOrange}{RGB}{239,126,45}
\definecolor{huangGold}{RGB}{196,145,15}
\definecolor{huangGreen}{RGB}{91,164,61}
\definecolor{huangGray}{RGB}{120,120,120}
\definecolor{huangPaleBlue}{RGB}{224,234,238}
\definecolor{huangPaleGreen}{RGB}{225,240,216}
\definecolor{huangPeach}{RGB}{246,185,139}
\definecolor{huangLavender}{RGB}{216,205,236}

\tikzset{
  huang arrow/.style={-Stealth, thick, draw=huangGray},
  huang dashed box/.style={draw, dashed, rounded corners=5pt, fill=huangPaleBlue, inner sep=5pt},
  huang bit/.style={draw, dashed, rounded corners=4pt, minimum width=0.42cm, minimum height=0.36cm, inner sep=1pt, font=\scriptsize\bfseries},
  huang gate/.style={draw, rounded corners=2pt, fill=huangPeach, minimum width=0.35cm, minimum height=0.54cm, inner sep=1pt},
  huang meas/.style={draw, rounded corners=2pt, fill=huangLavender, minimum width=0.24cm, minimum height=0.24cm, inner sep=0pt},
  huang label/.style={font=\scriptsize},
  huang title/.style={font=\bfseries}
}

\providecommand{\HuangAtom}[2]{%
  \begin{scope}[shift={(#1,#2)}, scale=0.18]
    \foreach \a in {0,60,...,300} {
      \draw[huangPurple, thick, rotate=\a] (0,0) ellipse (0.36 and 0.9);
    }
    \fill[huangGreen!85!black] (0,0) circle (0.11);
  \end{scope}%
}

\providecommand{\HuangLabPatch}[3][1]{%
  \begin{scope}[shift={(#2,#3)}, scale=#1]
    \node[huang dashed box, minimum width=2.25cm, minimum height=1.55cm] (labpatch) at (0,0) {};
    \foreach \x/\y in {-0.7/0.43,0/0.43,0.7/0.43,-0.7/-0.05,0/-0.05,0.7/-0.05,-0.35/-0.53,0.35/-0.53}
      {\HuangAtom{\x}{\y}}
  \end{scope}%
}

\providecommand{\HuangLabStack}[3][1]{%
  \begin{scope}[shift={(#2,#3)}, scale=#1]
    \foreach \dx in {-0.42,-0.18,0.06} {
      \node[huang dashed box, minimum width=1.75cm, minimum height=1.25cm, fill=huangPaleBlue!85] at (\dx,0) {};
    }
    \foreach \x/\y in {-0.58/0.36,0/0.36,0.58/0.36,-0.58/-0.03,0/-0.03,0.58/-0.03,-0.3/-0.43,0.3/-0.43}
      {\HuangAtom{\x}{\y}}
  \end{scope}%
}

\providecommand{\HuangMonitor}[3][1]{%
  \begin{scope}[shift={(#2,#3)}, scale=#1]
    \fill[black] (-0.72,-0.54) rectangle (0.72,0.55);
    \fill[huangLavender!60] (-0.6,-0.4) rectangle (0.6,0.4);
    \fill[white, opacity=0.72] (-0.55,-0.35) -- (-0.12,0.35) -- (0.6,0.35) -- (0.6,-0.35) -- cycle;
    \node[font=\scriptsize\bfseries, text=huangGray!80] at (0.12,0.02) {Model $\Psi$};
    \fill[black] (0,-0.47) circle (0.06);
    \fill[black] (-0.13,-0.8) rectangle (0.13,-0.54);
    \draw[black, line width=1.5pt] (-0.38,-0.84) -- (0.38,-0.84);
  \end{scope}%
}

\providecommand{\HuangMeasurementStack}[3][1]{%
  \begin{scope}[shift={(#2,#3)}, scale=#1]
    \foreach \i in {0,...,6} {
      \node[huang meas] at (0,0.18*\i) {};
      \draw[huangGold!80!black] (-0.08,0.18*\i-0.02) -- (0.08,0.18*\i+0.07);
    }
  \end{scope}%
}

\providecommand{\HuangCircuitGrid}[4][1]{%
  \begin{scope}[shift={(#2,#3)}, scale=#1]
    \foreach \x in {-0.75,-0.45,-0.15,0.15,0.45,0.75}
      \draw[huangGray] (\x,0.9) -- (\x,-1.05);
    \foreach \x/\y/\w/\h in {
      -0.75/0.55/.34/.45,-0.45/0.18/.24/.42,-0.15/0.55/.24/.52,0.15/0.18/.50/.42,0.45/0.55/.34/.45,0.75/0.18/.24/.42,
      -0.75/-0.25/.24/.46,-0.45/-0.63/.52/.42,-0.15/-0.25/.24/.52,0.15/-0.63/.52/.42,0.45/-0.25/.24/.46,0.75/-0.63/.34/.42}
      \node[huang gate, minimum width=\w cm, minimum height=\h cm] at (\x,\y) {};
    \node[font=\scriptsize] at (0,1.12) {$#4$};
  \end{scope}%
}

\providecommand{\HuangNeuralNet}[3][1]{%
  \begin{scope}[shift={(#2,#3)}, scale=#1]
    \node[draw, dashed, rounded corners=8pt, fill=huangPaleGreen, minimum width=1.65cm, minimum height=2.45cm] at (0,0) {};
    \foreach \x in {-0.45,0,0.45} \fill[huangBlue!38] (\x,0.75) circle (0.08);
    \foreach \x in {-0.58,-0.18,0.22,0.58} \fill[huangBlue!38] (\x,0.15) circle (0.08);
    \foreach \x in {-0.58,-0.18,0.22,0.58} \fill[huangBlue!38] (\x,-0.45) circle (0.08);
    \fill[huangBlue!38] (0,-1.02) circle (0.08);
    \foreach \a in {-0.45,0,0.45} \foreach \b in {-0.58,-0.18,0.22,0.58}
      \draw[huangGray!65, thin] (\a,0.75) -- (\b,0.15);
    \foreach \a in {-0.58,-0.18,0.22,0.58} \foreach \b in {-0.58,-0.18,0.22,0.58}
      \draw[huangGray!65, thin] (\a,0.15) -- (\b,-0.45);
    \foreach \a in {-0.58,-0.18,0.22,0.58}
      \draw[huangGray!65, thin] (\a,-0.45) -- (0,-1.02);
  \end{scope}%
}

\providecommand{\HuangDeviceIcon}[3][1]{%
  \begin{scope}[shift={(#2,#3)}, scale=#1]
    \fill[brown!80!black] (-0.72,0.82) rectangle (0.72,1.12);
    \fill[brown!70!black] (-0.82,0.72) rectangle (0.82,0.78);
    \foreach \x in {-0.48,-0.28,0.28,0.48}
      \draw[line width=1.5pt, huangGold] (\x,0.7) -- (\x,-0.55);
    \foreach \y in {0.42,0.12,-0.18}
      \draw[line width=1.2pt, brown!70!black] (-0.65,\y) -- (0.65,\y);
    \foreach \x in {-0.48,-0.38,0.38,0.48}
      \foreach \y in {0.28,-0.08}
        \draw[line width=1pt, huangGold] (\x,\y) circle (0.055);
    \foreach \x in {-0.52,-0.42,0.42,0.52}
      \foreach \y in {-0.38,-0.48,-0.58}
        \draw[line width=1.5pt, huangGold] (\x-0.08,\y) -- (\x+0.08,\y);
    \fill[huangGold!70] (-0.28,-0.6) rectangle (0.28,-0.85);
    \fill[brown!85!black] (-0.15,-0.85) rectangle (0.15,-1.18);
  \end{scope}%
}

\providecommand{\HuangHypercube}[3][1]{%
  \begin{scope}[shift={(#2,#3)}, scale=#1]
    \coordinate (z000) at (-1.2,-0.25);
    \coordinate (z001) at (-0.25,0.85);
    \coordinate (z010) at (-0.25,-0.25);
    \coordinate (z011) at (0.65,0.85);
    \coordinate (z100) at (-0.25,-1.2);
    \coordinate (z101) at (0.65,-0.25);
    \coordinate (z110) at (0.65,-1.2);
    \coordinate (z111) at (1.55,-0.25);
    \foreach \a/\b in {z000/z001,z000/z010,z000/z100,z001/z011,z001/z101,z010/z011,z010/z110,z100/z101,z100/z110,z011/z111,z101/z111,z110/z111}
      \draw[line width=1.5pt, huangGray] (\a) -- (\b);
    \foreach \p/\lab/\anch in {z000/000/below left,z001/001/above,z010/010/above left,z011/011/above,z100/100/below,z101/101/above,z110/110/below,z111/111/right}
      \fill[huangRed] (\p) circle (0.08) node[\anch, font=\scriptsize, text=black] {\lab};
  \end{scope}%
}
}
\begin{tikzpicture}[font=\footnotesize]

\node[font=\bfseries, align=center] at (-3.9,1.85) {Benchmarking\\quantum devices};
\HuangDeviceIcon[1.0]{-3.9}{0.05}

\draw[dashed, thick, huangGray!65] (-1.55,2.1) -- (-1.55,-1.8);

\node[font=\bfseries, align=center] at (0,1.85) {ML tomography\\of quantum states};
\node at (0,1.05) {$x\in\{0,1\}^n$};
\HuangNeuralNet[1.0]{0}{0}
\draw[huang arrow] (0,-1.1) -- (0,-1.55);
\node at (0,-1.75) {$\psi_\theta(x)$};

\draw[dashed, thick, huangGray!65] (1.55,2.1) -- (1.55,-1.8);

\node[font=\bfseries, align=center] at (3.9,1.85) {Optimizing\\quantum circuits};
\HuangCircuitGrid[1.0]{3.9}{0}{|0^n\rangle}
\node at (3.9,-1.35) {$|\psi\rangle$};

\end{tikzpicture}
}
    \end{center}

    \begin{itemize}
        \item Benchmark a noisy quantum device against its intended output.
        \item Train or certify neural quantum-state models using fidelity as the objective.
        \item Optimize parametrized circuits toward a target state without implementing the target projector.
    \end{itemize}
\end{frame}

\begin{frame}{What We Need From the Protocol}
    \begin{columns}[T,onlytextwidth]
        \begin{column}{0.5\textwidth}
            \begin{enumerate}
                \item Single-qubit measurements
                \item Only use classical description of target
            \end{enumerate}
        \end{column}
        \begin{column}{0.46\textwidth}
            \centering
            \resizebox{\linewidth}{!}{\IfFileExists{../diagrams/huang_templates.tex}{\definecolor{huangPurple}{RGB}{113,49,169}
\definecolor{huangBlue}{RGB}{65,109,190}
\definecolor{huangRed}{RGB}{196,0,0}
\definecolor{huangOrange}{RGB}{239,126,45}
\definecolor{huangGold}{RGB}{196,145,15}
\definecolor{huangGreen}{RGB}{91,164,61}
\definecolor{huangGray}{RGB}{120,120,120}
\definecolor{huangPaleBlue}{RGB}{224,234,238}
\definecolor{huangPaleGreen}{RGB}{225,240,216}
\definecolor{huangPeach}{RGB}{246,185,139}
\definecolor{huangLavender}{RGB}{216,205,236}

\tikzset{
  huang arrow/.style={-Stealth, thick, draw=huangGray},
  huang dashed box/.style={draw, dashed, rounded corners=5pt, fill=huangPaleBlue, inner sep=5pt},
  huang bit/.style={draw, dashed, rounded corners=4pt, minimum width=0.42cm, minimum height=0.36cm, inner sep=1pt, font=\scriptsize\bfseries},
  huang gate/.style={draw, rounded corners=2pt, fill=huangPeach, minimum width=0.35cm, minimum height=0.54cm, inner sep=1pt},
  huang meas/.style={draw, rounded corners=2pt, fill=huangLavender, minimum width=0.24cm, minimum height=0.24cm, inner sep=0pt},
  huang label/.style={font=\scriptsize},
  huang title/.style={font=\bfseries}
}

\providecommand{\HuangAtom}[2]{%
  \begin{scope}[shift={(#1,#2)}, scale=0.18]
    \foreach \a in {0,60,...,300} {
      \draw[huangPurple, thick, rotate=\a] (0,0) ellipse (0.36 and 0.9);
    }
    \fill[huangGreen!85!black] (0,0) circle (0.11);
  \end{scope}%
}

\providecommand{\HuangLabPatch}[3][1]{%
  \begin{scope}[shift={(#2,#3)}, scale=#1]
    \node[huang dashed box, minimum width=2.25cm, minimum height=1.55cm] (labpatch) at (0,0) {};
    \foreach \x/\y in {-0.7/0.43,0/0.43,0.7/0.43,-0.7/-0.05,0/-0.05,0.7/-0.05,-0.35/-0.53,0.35/-0.53}
      {\HuangAtom{\x}{\y}}
  \end{scope}%
}

\providecommand{\HuangLabStack}[3][1]{%
  \begin{scope}[shift={(#2,#3)}, scale=#1]
    \foreach \dx in {-0.42,-0.18,0.06} {
      \node[huang dashed box, minimum width=1.75cm, minimum height=1.25cm, fill=huangPaleBlue!85] at (\dx,0) {};
    }
    \foreach \x/\y in {-0.58/0.36,0/0.36,0.58/0.36,-0.58/-0.03,0/-0.03,0.58/-0.03,-0.3/-0.43,0.3/-0.43}
      {\HuangAtom{\x}{\y}}
  \end{scope}%
}

\providecommand{\HuangMonitor}[3][1]{%
  \begin{scope}[shift={(#2,#3)}, scale=#1]
    \fill[black] (-0.72,-0.54) rectangle (0.72,0.55);
    \fill[huangLavender!60] (-0.6,-0.4) rectangle (0.6,0.4);
    \fill[white, opacity=0.72] (-0.55,-0.35) -- (-0.12,0.35) -- (0.6,0.35) -- (0.6,-0.35) -- cycle;
    \node[font=\scriptsize\bfseries, text=huangGray!80] at (0.12,0.02) {Model $\Psi$};
    \fill[black] (0,-0.47) circle (0.06);
    \fill[black] (-0.13,-0.8) rectangle (0.13,-0.54);
    \draw[black, line width=1.5pt] (-0.38,-0.84) -- (0.38,-0.84);
  \end{scope}%
}

\providecommand{\HuangMeasurementStack}[3][1]{%
  \begin{scope}[shift={(#2,#3)}, scale=#1]
    \foreach \i in {0,...,6} {
      \node[huang meas] at (0,0.18*\i) {};
      \draw[huangGold!80!black] (-0.08,0.18*\i-0.02) -- (0.08,0.18*\i+0.07);
    }
  \end{scope}%
}

\providecommand{\HuangCircuitGrid}[4][1]{%
  \begin{scope}[shift={(#2,#3)}, scale=#1]
    \foreach \x in {-0.75,-0.45,-0.15,0.15,0.45,0.75}
      \draw[huangGray] (\x,0.9) -- (\x,-1.05);
    \foreach \x/\y/\w/\h in {
      -0.75/0.55/.34/.45,-0.45/0.18/.24/.42,-0.15/0.55/.24/.52,0.15/0.18/.50/.42,0.45/0.55/.34/.45,0.75/0.18/.24/.42,
      -0.75/-0.25/.24/.46,-0.45/-0.63/.52/.42,-0.15/-0.25/.24/.52,0.15/-0.63/.52/.42,0.45/-0.25/.24/.46,0.75/-0.63/.34/.42}
      \node[huang gate, minimum width=\w cm, minimum height=\h cm] at (\x,\y) {};
    \node[font=\scriptsize] at (0,1.12) {$#4$};
  \end{scope}%
}

\providecommand{\HuangNeuralNet}[3][1]{%
  \begin{scope}[shift={(#2,#3)}, scale=#1]
    \node[draw, dashed, rounded corners=8pt, fill=huangPaleGreen, minimum width=1.65cm, minimum height=2.45cm] at (0,0) {};
    \foreach \x in {-0.45,0,0.45} \fill[huangBlue!38] (\x,0.75) circle (0.08);
    \foreach \x in {-0.58,-0.18,0.22,0.58} \fill[huangBlue!38] (\x,0.15) circle (0.08);
    \foreach \x in {-0.58,-0.18,0.22,0.58} \fill[huangBlue!38] (\x,-0.45) circle (0.08);
    \fill[huangBlue!38] (0,-1.02) circle (0.08);
    \foreach \a in {-0.45,0,0.45} \foreach \b in {-0.58,-0.18,0.22,0.58}
      \draw[huangGray!65, thin] (\a,0.75) -- (\b,0.15);
    \foreach \a in {-0.58,-0.18,0.22,0.58} \foreach \b in {-0.58,-0.18,0.22,0.58}
      \draw[huangGray!65, thin] (\a,0.15) -- (\b,-0.45);
    \foreach \a in {-0.58,-0.18,0.22,0.58}
      \draw[huangGray!65, thin] (\a,-0.45) -- (0,-1.02);
  \end{scope}%
}

\providecommand{\HuangDeviceIcon}[3][1]{%
  \begin{scope}[shift={(#2,#3)}, scale=#1]
    \fill[brown!80!black] (-0.72,0.82) rectangle (0.72,1.12);
    \fill[brown!70!black] (-0.82,0.72) rectangle (0.82,0.78);
    \foreach \x in {-0.48,-0.28,0.28,0.48}
      \draw[line width=1.5pt, huangGold] (\x,0.7) -- (\x,-0.55);
    \foreach \y in {0.42,0.12,-0.18}
      \draw[line width=1.2pt, brown!70!black] (-0.65,\y) -- (0.65,\y);
    \foreach \x in {-0.48,-0.38,0.38,0.48}
      \foreach \y in {0.28,-0.08}
        \draw[line width=1pt, huangGold] (\x,\y) circle (0.055);
    \foreach \x in {-0.52,-0.42,0.42,0.52}
      \foreach \y in {-0.38,-0.48,-0.58}
        \draw[line width=1.5pt, huangGold] (\x-0.08,\y) -- (\x+0.08,\y);
    \fill[huangGold!70] (-0.28,-0.6) rectangle (0.28,-0.85);
    \fill[brown!85!black] (-0.15,-0.85) rectangle (0.15,-1.18);
  \end{scope}%
}

\providecommand{\HuangHypercube}[3][1]{%
  \begin{scope}[shift={(#2,#3)}, scale=#1]
    \coordinate (z000) at (-1.2,-0.25);
    \coordinate (z001) at (-0.25,0.85);
    \coordinate (z010) at (-0.25,-0.25);
    \coordinate (z011) at (0.65,0.85);
    \coordinate (z100) at (-0.25,-1.2);
    \coordinate (z101) at (0.65,-0.25);
    \coordinate (z110) at (0.65,-1.2);
    \coordinate (z111) at (1.55,-0.25);
    \foreach \a/\b in {z000/z001,z000/z010,z000/z100,z001/z011,z001/z101,z010/z011,z010/z110,z100/z101,z100/z110,z011/z111,z101/z111,z110/z111}
      \draw[line width=1.5pt, huangGray] (\a) -- (\b);
    \foreach \p/\lab/\anch in {z000/000/below left,z001/001/above,z010/010/above left,z011/011/above,z100/100/below,z101/101/above,z110/110/below,z111/111/right}
      \fill[huangRed] (\p) circle (0.08) node[\anch, font=\scriptsize, text=black] {\lab};
  \end{scope}%
}
}{\definecolor{huangPurple}{RGB}{113,49,169}
\definecolor{huangBlue}{RGB}{65,109,190}
\definecolor{huangRed}{RGB}{196,0,0}
\definecolor{huangOrange}{RGB}{239,126,45}
\definecolor{huangGold}{RGB}{196,145,15}
\definecolor{huangGreen}{RGB}{91,164,61}
\definecolor{huangGray}{RGB}{120,120,120}
\definecolor{huangPaleBlue}{RGB}{224,234,238}
\definecolor{huangPaleGreen}{RGB}{225,240,216}
\definecolor{huangPeach}{RGB}{246,185,139}
\definecolor{huangLavender}{RGB}{216,205,236}

\tikzset{
  huang arrow/.style={-Stealth, thick, draw=huangGray},
  huang dashed box/.style={draw, dashed, rounded corners=5pt, fill=huangPaleBlue, inner sep=5pt},
  huang bit/.style={draw, dashed, rounded corners=4pt, minimum width=0.42cm, minimum height=0.36cm, inner sep=1pt, font=\scriptsize\bfseries},
  huang gate/.style={draw, rounded corners=2pt, fill=huangPeach, minimum width=0.35cm, minimum height=0.54cm, inner sep=1pt},
  huang meas/.style={draw, rounded corners=2pt, fill=huangLavender, minimum width=0.24cm, minimum height=0.24cm, inner sep=0pt},
  huang label/.style={font=\scriptsize},
  huang title/.style={font=\bfseries}
}

\providecommand{\HuangAtom}[2]{%
  \begin{scope}[shift={(#1,#2)}, scale=0.18]
    \foreach \a in {0,60,...,300} {
      \draw[huangPurple, thick, rotate=\a] (0,0) ellipse (0.36 and 0.9);
    }
    \fill[huangGreen!85!black] (0,0) circle (0.11);
  \end{scope}%
}

\providecommand{\HuangLabPatch}[3][1]{%
  \begin{scope}[shift={(#2,#3)}, scale=#1]
    \node[huang dashed box, minimum width=2.25cm, minimum height=1.55cm] (labpatch) at (0,0) {};
    \foreach \x/\y in {-0.7/0.43,0/0.43,0.7/0.43,-0.7/-0.05,0/-0.05,0.7/-0.05,-0.35/-0.53,0.35/-0.53}
      {\HuangAtom{\x}{\y}}
  \end{scope}%
}

\providecommand{\HuangLabStack}[3][1]{%
  \begin{scope}[shift={(#2,#3)}, scale=#1]
    \foreach \dx in {-0.42,-0.18,0.06} {
      \node[huang dashed box, minimum width=1.75cm, minimum height=1.25cm, fill=huangPaleBlue!85] at (\dx,0) {};
    }
    \foreach \x/\y in {-0.58/0.36,0/0.36,0.58/0.36,-0.58/-0.03,0/-0.03,0.58/-0.03,-0.3/-0.43,0.3/-0.43}
      {\HuangAtom{\x}{\y}}
  \end{scope}%
}

\providecommand{\HuangMonitor}[3][1]{%
  \begin{scope}[shift={(#2,#3)}, scale=#1]
    \fill[black] (-0.72,-0.54) rectangle (0.72,0.55);
    \fill[huangLavender!60] (-0.6,-0.4) rectangle (0.6,0.4);
    \fill[white, opacity=0.72] (-0.55,-0.35) -- (-0.12,0.35) -- (0.6,0.35) -- (0.6,-0.35) -- cycle;
    \node[font=\scriptsize\bfseries, text=huangGray!80] at (0.12,0.02) {Model $\Psi$};
    \fill[black] (0,-0.47) circle (0.06);
    \fill[black] (-0.13,-0.8) rectangle (0.13,-0.54);
    \draw[black, line width=1.5pt] (-0.38,-0.84) -- (0.38,-0.84);
  \end{scope}%
}

\providecommand{\HuangMeasurementStack}[3][1]{%
  \begin{scope}[shift={(#2,#3)}, scale=#1]
    \foreach \i in {0,...,6} {
      \node[huang meas] at (0,0.18*\i) {};
      \draw[huangGold!80!black] (-0.08,0.18*\i-0.02) -- (0.08,0.18*\i+0.07);
    }
  \end{scope}%
}

\providecommand{\HuangCircuitGrid}[4][1]{%
  \begin{scope}[shift={(#2,#3)}, scale=#1]
    \foreach \x in {-0.75,-0.45,-0.15,0.15,0.45,0.75}
      \draw[huangGray] (\x,0.9) -- (\x,-1.05);
    \foreach \x/\y/\w/\h in {
      -0.75/0.55/.34/.45,-0.45/0.18/.24/.42,-0.15/0.55/.24/.52,0.15/0.18/.50/.42,0.45/0.55/.34/.45,0.75/0.18/.24/.42,
      -0.75/-0.25/.24/.46,-0.45/-0.63/.52/.42,-0.15/-0.25/.24/.52,0.15/-0.63/.52/.42,0.45/-0.25/.24/.46,0.75/-0.63/.34/.42}
      \node[huang gate, minimum width=\w cm, minimum height=\h cm] at (\x,\y) {};
    \node[font=\scriptsize] at (0,1.12) {$#4$};
  \end{scope}%
}

\providecommand{\HuangNeuralNet}[3][1]{%
  \begin{scope}[shift={(#2,#3)}, scale=#1]
    \node[draw, dashed, rounded corners=8pt, fill=huangPaleGreen, minimum width=1.65cm, minimum height=2.45cm] at (0,0) {};
    \foreach \x in {-0.45,0,0.45} \fill[huangBlue!38] (\x,0.75) circle (0.08);
    \foreach \x in {-0.58,-0.18,0.22,0.58} \fill[huangBlue!38] (\x,0.15) circle (0.08);
    \foreach \x in {-0.58,-0.18,0.22,0.58} \fill[huangBlue!38] (\x,-0.45) circle (0.08);
    \fill[huangBlue!38] (0,-1.02) circle (0.08);
    \foreach \a in {-0.45,0,0.45} \foreach \b in {-0.58,-0.18,0.22,0.58}
      \draw[huangGray!65, thin] (\a,0.75) -- (\b,0.15);
    \foreach \a in {-0.58,-0.18,0.22,0.58} \foreach \b in {-0.58,-0.18,0.22,0.58}
      \draw[huangGray!65, thin] (\a,0.15) -- (\b,-0.45);
    \foreach \a in {-0.58,-0.18,0.22,0.58}
      \draw[huangGray!65, thin] (\a,-0.45) -- (0,-1.02);
  \end{scope}%
}

\providecommand{\HuangDeviceIcon}[3][1]{%
  \begin{scope}[shift={(#2,#3)}, scale=#1]
    \fill[brown!80!black] (-0.72,0.82) rectangle (0.72,1.12);
    \fill[brown!70!black] (-0.82,0.72) rectangle (0.82,0.78);
    \foreach \x in {-0.48,-0.28,0.28,0.48}
      \draw[line width=1.5pt, huangGold] (\x,0.7) -- (\x,-0.55);
    \foreach \y in {0.42,0.12,-0.18}
      \draw[line width=1.2pt, brown!70!black] (-0.65,\y) -- (0.65,\y);
    \foreach \x in {-0.48,-0.38,0.38,0.48}
      \foreach \y in {0.28,-0.08}
        \draw[line width=1pt, huangGold] (\x,\y) circle (0.055);
    \foreach \x in {-0.52,-0.42,0.42,0.52}
      \foreach \y in {-0.38,-0.48,-0.58}
        \draw[line width=1.5pt, huangGold] (\x-0.08,\y) -- (\x+0.08,\y);
    \fill[huangGold!70] (-0.28,-0.6) rectangle (0.28,-0.85);
    \fill[brown!85!black] (-0.15,-0.85) rectangle (0.15,-1.18);
  \end{scope}%
}

\providecommand{\HuangHypercube}[3][1]{%
  \begin{scope}[shift={(#2,#3)}, scale=#1]
    \coordinate (z000) at (-1.2,-0.25);
    \coordinate (z001) at (-0.25,0.85);
    \coordinate (z010) at (-0.25,-0.25);
    \coordinate (z011) at (0.65,0.85);
    \coordinate (z100) at (-0.25,-1.2);
    \coordinate (z101) at (0.65,-0.25);
    \coordinate (z110) at (0.65,-1.2);
    \coordinate (z111) at (1.55,-0.25);
    \foreach \a/\b in {z000/z001,z000/z010,z000/z100,z001/z011,z001/z101,z010/z011,z010/z110,z100/z101,z100/z110,z011/z111,z101/z111,z110/z111}
      \draw[line width=1.5pt, huangGray] (\a) -- (\b);
    \foreach \p/\lab/\anch in {z000/000/below left,z001/001/above,z010/010/above left,z011/011/above,z100/100/below,z101/101/above,z110/110/below,z111/111/right}
      \fill[huangRed] (\p) circle (0.08) node[\anch, font=\scriptsize, text=black] {\lab};
  \end{scope}%
}
}
\begin{tikzpicture}[>=Stealth, font=\footnotesize]

\node[font=\bfseries] at (-2.2,1.55) {Lab state $\rho$};
\node[draw, rounded corners=12pt, fill=huangPaleBlue, minimum width=2.15cm, minimum height=1.8cm] at (-2.2,0.55) {};
\foreach \x/\y in {-2.55/1.0,-1.9/1.0,-2.85/0.55,-2.2/0.55,-2.55/0.1,-1.9/0.1}
  \HuangAtom{\x}{\y}
\node[huang bit, fill=orange!18] at (-1.45,0.55) {};
\HuangAtom{-1.45}{0.55}
\node[anchor=west, font=\itshape\bfseries, text=huangGray] at (-1.2,0.75) {qubit $k$};

\node[font=\bfseries] at (2.5,1.55) {Target state};
\node at (2.5,1.22) {$|\psi\rangle=\sum_x\psi(x)|x\rangle$};
\HuangMonitor[0.95]{2.5}{0.35}

\node[font=\large] at (0,-0.7) {Naive idea:};
\node[font=\large] at (0,-1.18) {Choose a qubit at random and compare locally};
\node[font=\large] at (0,-1.75) {Issue: one-qubit marginal $\approx$ maximally mixed};

\end{tikzpicture}
}
        \end{column}
    \end{columns}

    \vspace{0.4em}
    \begin{center}
        The algorithm resolves the ``locally maximally mixed'' obstacle by conditioning on previous measurement outcomes.
    \end{center}
\end{frame}
